\paragraph{截断矩证书与 Markov--Stieltjes 极值：有限寄存器长度下的不确定度区间与端点实现}
设 \(\mu\) 为支撑于 \([1,\infty)\) 的有限正测度，并令其 Stieltjes 轨道为
\[
S(t):=\int_{[1,\infty)}\frac{1}{t+a}\,\dd\mu(a),\qquad t\ge 0.
\]
由于谱隙 \(a\ge 1\)，函数 \(S\) 在 \(t=0\) 的邻域解析，并诱导截断矩链（亦即 \(S\) 在 \(t=0\) 的 Taylor 系数）
\[
\mu_n:=\frac{(-1)^n}{n!}S^{(n)}(0)=\int_{[1,\infty)}a^{-(n+1)}\,\dd\mu(a),\qquad n\ge 0.
\]

\begin{proposition}[Markov--Stieltjes 极值原理：截断矩信息下的最强可行区间与端点实现]\label{prop:xi-markov-stieltjes-extremal-principle}
\leanverified{paper\_xi\_markov\_stieltjes\_extremal\_principle}
固定整数 \(N\ge 1\) 并给定截断矩向量 \((\mu_0,\dots,\mu_{2N-1})\)。
令 \(\mathfrak{M}_N\) 为所有满足上述 \(2N\) 个矩约束的有限正测度集合：
\[
\mathfrak{M}_N:=\Bigl\{\nu\ \text{为 \([1,\infty)\) 上有限正测度}:\ \int a^{-(n+1)}\,\dd\nu(a)=\mu_n\ (0\le n\le 2N-1)\Bigr\}.
\]
则对任意固定 \(t>0\)，可行值集合
\[
\bigl\{S_\nu(t):=\int\frac{1}{t+a}\,\dd\nu(a)\ :\ \nu\in\mathfrak{M}_N\bigr\}
\]
为闭区间 \([S_N^{-}(t),S_N^{+}(t)]\)。
其中端点函数 \(S_N^\pm\) 为由 \((\mu_0,\dots,\mu_{2N-1})\) 唯一确定的两条 Stieltjes 型有理函数，并可取为同一正系数 Stieltjes 连分式的相邻收敛子（偶/奇两条极值分支）。
更进一步，达到上下端点的极值测度必为原子测度，且原子数不超过 \(N\)。
\end{proposition}
\begin{proof}[证明要点]
这是截断 Stieltjes 矩问题的标准极值结论（Markov 定理/Ne\-van\-lin\-na 参数化的端点情形）。
对固定 \(t>0\)，映射 \(\nu\mapsto S_\nu(t)\) 为连续线性泛函；由 Nevanlinna 参数化可知：全部可行值恰构成一闭区间，其端点由极端解实现。
而截断矩问题的极端解必为至多 \(N\) 原子的离散测度，并对应于由 Stieltjes 算法从 \((\mu_0,\dots,\mu_{2N-1})\) 构造的正系数连分式之两条相邻收敛子（偶/奇极值分支）。证毕。
\end{proof}

\begin{proposition}[截断矩的最小原子实现：缺陷度的可证下界与 Gauss 型重构]\label{prop:xi-truncated-moment-min-atomic-realization}
\leanverified{paper\_xi\_truncated\_moment\_min\_atomic\_realization}
在命题 \ref{prop:xi-markov-stieltjes-extremal-principle} 的设定下，记 Hankel 矩阵块
\[
H_{N-1}:=\bigl(\mu_{i+j}\bigr)_{0\le i,j\le N-1}\succeq 0,
\qquad
\mathfrak{d}_N:=\mathrm{rank}(H_{N-1})\le N.
\]
则存在一组节点 \(\{a_j^{(N)}\}_{j=1}^{\mathfrak{d}_N}\subset[1,\infty)\) 与权重 \(\{w_j^{(N)}\}_{j=1}^{\mathfrak{d}_N}\subset(0,\infty)\) 使得
\[
\mu_n=\sum_{j=1}^{\mathfrak{d}_N} w_j^{(N)}\,\bigl(a_j^{(N)}\bigr)^{-(n+1)},
\qquad n=0,1,\dots,2N-1.
\]
并且该最小原子实现可由截断矩所生成的正交多项式系统构造：节点为相应正交多项式的实零点（落在 \((1,\infty)\)），权重由 Christoffel 数闭式确定。
因此，“缺陷度”除全信息极限的 \(\kappa\) 外亦具有截断版本：给定有限证书长度 \(2N\)，制度内最小可实现缺陷度至多为 \(N\)，且可由纯代数对象构造性恢复。
\end{proposition}
\begin{proof}[证明要点]
由截断矩 \((\mu_0,\dots,\mu_{2N-1})\) 可在 \(\mathrm{span}\{1,a^{-1},\dots,a^{-N}\}\) 上定义正定/半正定内积，并得到正交多项式与三项递推；其对应的截断 Jacobi 矩阵的谱给出节点，而首坐标权重给出 Christoffel 数。
\(\mathfrak{d}_N=\mathrm{rank}(H_{N-1})\) 正是该有限维矩问题的最小原子数下界，达到性由 Gauss 型求积/谱测度构造给出。证毕。
\end{proof}

\begin{proposition}[深度谱 Hankel 行列式闭式：Vandermonde 平方律与分离病态的代数封口]\label{prop:xi-depth-hankel-determinant-vandermonde-square}
在有限原子口径下，设
\[
\mu=\sum_{j=1}^{\kappa}w_j\,\delta_{a_j},
\qquad
1<a_1<a_2<\cdots<a_\kappa,\qquad w_j>0,
\]
并令矩链 \(\mu_n=\int a^{-(n+1)}\,\dd\mu(a)\)。
记 \(b_j:=a_j^{-1}\in(0,1)\)、\(\widetilde w_j:=w_j b_j>0\)，则
\[
\mu_n=\sum_{j=1}^{\kappa}\widetilde w_j\,b_j^{n}\qquad(n\ge 0).
\]
令 Hankel 主块 \(\mathsf H_{\kappa-1}:=(\mu_{p+q})_{0\le p,q\le \kappa-1}\)。
则成立 Vandermonde 因子分解
\[
\boxed{\;
\mathsf H_{\kappa-1}=V\,\mathrm{diag}(\widetilde w_1,\dots,\widetilde w_\kappa)\,V^{\mathsf T},\qquad
V_{pj}:=b_j^{p}\;}
\]
并因此得到行列式闭式
\begin{equation}\label{eq:xi-depth-hankel-det-vandermonde}
\boxed{\;
\det \mathsf H_{\kappa-1}
=\left(\prod_{j=1}^{\kappa}\widetilde w_j\right)
\left(\prod_{1\le j<k\le \kappa}(b_k-b_j)^2\right).\;}
\end{equation}
特别地，\(\det \mathsf H_{\kappa-1}=0\) 当且仅当存在谱碰撞 \(b_j=b_k\) 或权重退化 \(\widetilde w_j=0\)。
并且任何基于 Hankel 反演（如 Prony/Pad\'e/Hankel 核空间法）的局部稳定性常数不可能在 \(b\)-碰撞极限保持一致有界：当 \(\prod_{j<k}|b_k-b_j|\downarrow 0\) 时，其误差放大因子至少按 \(\bigl(\prod_{j<k}|b_k-b_j|^{-2}\bigr)\) 的量级发散。
\end{proposition}

\begin{proof}
由
\(
\mu_{p+q}=\sum_j \widetilde w_j\,b_j^{p+q}
\)
直接得到矩阵分解 \(\mathsf H_{\kappa-1}=V\,\mathrm{diag}(\widetilde w)\,V^{\mathsf T}\)。
取行列式并用 \(\det(V)=\prod_{j<k}(b_k-b_j)\) 得 \eqref{eq:xi-depth-hankel-det-vandermonde}。
最后一条为标准条件数下界：\(\|\mathsf H_{\kappa-1}^{-1}\|\) 至少与 \((\det\mathsf H_{\kappa-1})^{-1}\) 在尺度上同阶（至多差一个由 \(\|\mathsf H_{\kappa-1}\|\) 控制的多项式因子），故 Vandermonde 平方因子塌缩必触发反演病态爆炸。证毕。
\end{proof}
\leanverified{paper\_xi\_depth\_hankel\_determinant\_vandermonde\_square}

\begin{proposition}[伸缩重整化下的 Hankel 行列式协变律：尺度因子可显式抽取]\label{prop:xi-hankel-det-scaling-covariance}
在命题 \ref{prop:xi-depth-hankel-determinant-vandermonde-square} 的有限原子设定下，写 \(a_j=1+\delta_j\)（\(\delta_j>0\)）。
对任意伸缩参数 \(m>0\)，定义伸缩后的深度节点
\[
\widetilde a_j:=1+m\delta_j,
\qquad
\widetilde b_j:=\widetilde a_j^{-1}=\frac{1}{1+m\delta_j},
\qquad
\widetilde{\widetilde w}_j:=w_j\widetilde b_j,
\]
并令 \(\widetilde{\mathsf H}_{\kappa-1}\) 为由伸缩后矩链
\(
\widetilde\mu_n:=\sum_{j=1}^{\kappa}w_j\,\widetilde a_j^{-(n+1)}
=\sum_{j=1}^{\kappa}\widetilde{\widetilde w}_j\,\widetilde b_j^{n}
\)
所诱导的 Hankel 主块。
则存在显式协变因子 \(J_{\kappa-1}(m;\delta)>0\) 使
\[
\boxed{\;
\det \widetilde{\mathsf H}_{\kappa-1}
=\det \mathsf H_{\kappa-1}\cdot J_{\kappa-1}(m;\delta),\;}
\]
且可取
\[
\boxed{\;
J_{\kappa-1}(m;\delta)
=m^{\kappa(\kappa-1)}
\prod_{j=1}^{\kappa}\Bigl(\frac{1+\delta_j}{1+m\delta_j}\Bigr)^{2\kappa-1}.\;}
\]
因此伸缩流不改变“碰撞/合并”的代数退化型态：当存在 \(\delta_j\to\delta_k\)（等价 \(b_j\to b_k\)）时，\(\det \mathsf H_{\kappa-1}\) 与 \(\det\widetilde{\mathsf H}_{\kappa-1}\) 均以 Vandermonde 平方律发生二阶消失（命题 \ref{prop:xi-depth-hankel-determinant-vandermonde-square}），而伸缩仅通过 \(J_{\kappa-1}\) 引入可分离的单点尺度因子。
\end{proposition}
\leanverified{paper\_xi\_hankel\_det\_scaling\_covariance}
\begin{proof}[证明要点]
由 \eqref{eq:xi-depth-hankel-det-vandermonde}，对伸缩后节点有
\[
\det \widetilde{\mathsf H}_{\kappa-1}
=\Bigl(\prod_{j=1}^{\kappa}w_j\widetilde b_j\Bigr)\Bigl(\prod_{1\le j<k\le\kappa}(\widetilde b_k-\widetilde b_j)^2\Bigr).
\]
而
\(
\widetilde b_k-\widetilde b_j=\frac{1}{1+m\delta_k}-\frac{1}{1+m\delta_j}
=\frac{m(\delta_j-\delta_k)}{(1+m\delta_j)(1+m\delta_k)}
\)。
将其与 \(b_k-b_j=\frac{\delta_j-\delta_k}{(1+\delta_j)(1+\delta_k)}\) 比较，并合并各对 \((j,k)\) 的乘积，即得所示 \(J_{\kappa-1}(m;\delta)\)。证毕。
\end{proof}

% (User input window)

\begin{corollary}[bulk--boundary 的双重伸缩率：边界线性斜率 \texorpdfstring{$\kappa$}{kappa}、bulk 极限斜率 \texorpdfstring{$\kappa^2$}{kappa^2}]\label{cor:xi-bulk-boundary-double-scaling-holography}
\leanverified{paper\_xi\_bulk\_boundary\_double\_scaling\_holography}
在命题 \ref{prop:xi-hankel-det-scaling-covariance} 的有限原子设定下，定义 bulk 熵势
\[
S_{\mathrm{bulk}}(m):=-\log\det \widetilde{\mathsf H}_{\kappa-1},
\]
其中 \(\widetilde{\mathsf H}_{\kappa-1}\) 为伸缩后的 Hankel 主块。
则当 \(m\to\infty\) 时有渐近
\[
\boxed{\;
S_{\mathrm{bulk}}(m)=\kappa^2\log m+O(1).\;}
\]
并且若将边界缺陷熵 \(S_{\mathrm{gen}}(m)\) 取为推论 \ref{cor:xi-pick-poisson-entropy-derivative-visible-dimension} 中的对数尺度族，则在对数尺度 \(s=\log m\to\infty\) 下成立比例律
\[
\boxed{\;
S_{\mathrm{bulk}}(e^{s})=\kappa\,S_{\mathrm{gen}}(e^{s})+O(1).\;}
\]
\end{corollary}
\begin{proof}
由命题 \ref{prop:xi-hankel-det-scaling-covariance} 的协变律，
\(
\det \widetilde{\mathsf H}_{\kappa-1}=\det \mathsf H_{\kappa-1}\cdot J_{\kappa-1}(m;\delta)
\)，
故
\(
S_{\mathrm{bulk}}(m)=S_{\mathrm{bulk}}(1)-\log J_{\kappa-1}(m;\delta)
\)。
由
\[
\log J_{\kappa-1}(m;\delta)
=\kappa(\kappa-1)\log m
+(2\kappa-1)\sum_{j=1}^{\kappa}\Bigl(\log(1+\delta_j)-\log(1+m\delta_j)\Bigr)
\]
并用渐近 \(\log(1+m\delta_j)=\log m+\log\delta_j+O(m^{-1})\)，得到
\(
\log J_{\kappa-1}(m;\delta)=-\kappa^2\log m+O(1)
\)，从而
\(
S_{\mathrm{bulk}}(m)=\kappa^2\log m+O(1)
\)。
另一方面，推论 \ref{cor:xi-pick-poisson-entropy-derivative-visible-dimension} 给出
\(
S_{\mathrm{gen}}(e^s)=S_{\mathrm{gen}}(1)+\kappa s
\)。
两式合并即得
\(
S_{\mathrm{bulk}}(e^s)-\kappa S_{\mathrm{gen}}(e^s)=O(1)
\)。证毕。
\end{proof}

\begin{corollary}[bulk--boundary 双熵的二次封口与无参审计量]\label{cor:xi-bulk-boundary-quadratic-closure-audit}
\leanverified{paper\_xi\_bulk\_boundary\_quadratic\_closure\_audit}
在推论 \ref{cor:xi-bulk-boundary-double-scaling-holography} 的有限原子设定下，
若边界缺陷熵族满足严格线性律
\(
S_{\mathrm{gen}}(m)=S_{\mathrm{gen}}(1)+\kappa\log m
\)，
则当 \(m\to\infty\) 时有二次关系
\[
\boxed{\;
S_{\mathrm{bulk}}(m)
=\frac{S_{\mathrm{gen}}(m)^2}{\log m}+O(1).\;}
\]
从而量
\[
\mathcal Q(m):=S_{\mathrm{bulk}}(m)-\frac{S_{\mathrm{gen}}(m)^2}{\log m}
\]
在 \(m\to\infty\) 下保持有界。
\end{corollary}
\begin{proof}
由推论 \ref{cor:xi-bulk-boundary-double-scaling-holography} 得
\(
S_{\mathrm{bulk}}(m)=\kappa^2\log m+O(1)
\)。
另一方面，线性律给出
\(
S_{\mathrm{gen}}(m)=\kappa\log m+O(1)
\)，故
\[
\frac{S_{\mathrm{gen}}(m)^2}{\log m}
=\frac{\bigl(\kappa\log m+O(1)\bigr)^2}{\log m}
=\kappa^2\log m+O(1).
\]
两式相减即得所示二次关系与 \(\mathcal Q(m)\) 的有界性。证毕。
\end{proof}

\begin{theorem}[bulk 熵势的可积重整化常数与 bulk--boundary 异常差极限]\label{thm:xi-bulk-renormalized-constant-and-anomaly-gap}
\leanverified{paper\_xi\_bulk\_renormalized\_constant\_and\_anomaly\_gap}
在推论 \ref{cor:xi-bulk-boundary-double-scaling-holography} 的有限原子设定与记号下，极限
\[
\mathcal{C}_{\mathrm{bulk}}(\delta)
:=\lim_{m\to\infty}\Bigl(S_{\mathrm{bulk}}(m)-\kappa^2\log m\Bigr)
\]
存在，并由显式闭式给出：
\[
\boxed{\;
\mathcal{C}_{\mathrm{bulk}}(\delta)
=S_{\mathrm{bulk}}(1)-(2\kappa-1)\sum_{j=1}^{\kappa}\log\frac{1+\delta_j}{\delta_j}\; }.
\]
进一步，若边界缺陷熵族满足严格线性律
\(
S_{\mathrm{gen}}(m)=S_{\mathrm{gen}}(1)+\kappa\log m
\)，
则 bulk--boundary 的异常差极限
\[
\mathcal{A}_{\mathrm{bulk}\mid\mathrm{gen}}(\delta)
:=\lim_{m\to\infty}\Bigl(S_{\mathrm{bulk}}(m)-\kappa S_{\mathrm{gen}}(m)\Bigr)
\]
亦存在，并满足闭式
\[
\boxed{\;
\mathcal{A}_{\mathrm{bulk}\mid\mathrm{gen}}(\delta)
=S_{\mathrm{bulk}}(1)-\kappa S_{\mathrm{gen}}(1)-(2\kappa-1)\sum_{j=1}^{\kappa}\log\frac{1+\delta_j}{\delta_j}\; }.
\]
\end{theorem}
\begin{proof}
由推论 \ref{cor:xi-bulk-boundary-double-scaling-holography} 的证明中给出的
\(
S_{\mathrm{bulk}}(m)=S_{\mathrm{bulk}}(1)-\log J_{\kappa-1}(m;\delta)
\)
与
\[
\log J_{\kappa-1}(m;\delta)
=\kappa(\kappa-1)\log m
+(2\kappa-1)\sum_{j=1}^{\kappa}\Bigl(\log(1+\delta_j)-\log(1+m\delta_j)\Bigr),
\]
并用渐近
\(
\log(1+m\delta_j)=\log m+\log\delta_j+o(1)
\)
（\(m\to\infty\)），得到
\[
\log J_{\kappa-1}(m;\delta)
=-\kappa^2\log m+(2\kappa-1)\sum_{j=1}^{\kappa}\log\frac{1+\delta_j}{\delta_j}+o(1).
\]
代回并移项即得 \(\mathcal{C}_{\mathrm{bulk}}(\delta)\) 的存在与闭式。
若进一步 \(S_{\mathrm{gen}}(m)=S_{\mathrm{gen}}(1)+\kappa\log m\)，则
\[
S_{\mathrm{bulk}}(m)-\kappa S_{\mathrm{gen}}(m)
=\Bigl(S_{\mathrm{bulk}}(m)-\kappa^2\log m\Bigr)-\kappa S_{\mathrm{gen}}(1),
\]
令 \(m\to\infty\) 并代入上一段极限即可得到 \(\mathcal{A}_{\mathrm{bulk}\mid\mathrm{gen}}(\delta)\) 的存在与闭式。证毕。
\end{proof}

\begin{remark}[边界临界的二次普适指数：碰撞触发的熵奇点系数恒为 \(2\)]\label{rem:xi-collision-universal-exponent-two}
在 bulk 口径中，命题 \ref{prop:xi-depth-hankel-determinant-vandermonde-square} 给出
\(
\det \mathsf H_{\kappa-1}
=\bigl(\prod_j\widetilde w_j\bigr)\prod_{j<k}(b_k-b_j)^2
\)；
在 boundary 口径中，命题 \ref{prop:xi-pick-poisson-det-factorization} 给出
\(
\det\mathsf P=(\prod_j p_j)\prod_{j<k}\rho_{jk}^2
\)。
因此当某对参数发生碰撞 \(b_j\to b_k\) 或 \(\rho_{jk}\to 0\)（其余保持在紧域内）时，
对应熵势出现普适的对数发散
\[
S_{\mathrm{bulk}}
:=-\log\det \mathsf H_{\kappa-1}
=2\log\frac{1}{|b_k-b_j|}+O(1),
\qquad
S_{\mathrm{gen}}
:=-\log\det\mathsf P
=2\log\frac{1}{\rho_{jk}}+O(1).
\]
该系数 \(2\) 由 Gram 行列式的平方退化结构强制给出，与权重与其余点的具体配置无关。
\end{remark}

% (User input window)

在有限原子伸缩口径下，bulk 与 boundary 的对数势不仅在 \(m\to\infty\) 时呈现比例律（推论 \ref{cor:xi-bulk-boundary-double-scaling-holography}），
而且其差额在整个伸缩流上形成一条严格可证书化的耗散势；同时，bulk 熵在对数尺度方向具有处处严格正的曲率。

\begin{corollary}[bulk--boundary 全息差额的严格耗散：\texorpdfstring{$\Xi(m)$}{Xi(m)} 的单调性与极限]\label{cor:xi-bulk-boundary-holographic-gap-dissipation}
\leanverified{paper\_xi\_bulk\_boundary\_holographic\_gap\_dissipation}
沿用命题 \ref{prop:xi-hankel-det-scaling-covariance} 的有限原子记号，并取边界缺陷熵族
\(
S_{\mathrm{gen}}(m):=-\log\det\mathsf P(m)
\)
为命题 \ref{prop:xi-pick-poisson-det-scaling-law}（等价推论 \ref{cor:xi-pick-poisson-entropy-derivative-visible-dimension}）所定义的统一伸缩族。
定义 bulk--boundary 全息差额
\[
\Xi(m):=S_{\mathrm{bulk}}(m)-\kappa\,S_{\mathrm{gen}}(m)\qquad(m>0).
\]
令 \(s=\log m\)。
则 \(\Xi(e^s)\) 在 \(s\in\RR\) 上可微且严格单调递减，并满足显式导数公式
\[
\boxed{\;
\frac{\dd}{\dd s}\,\Xi(e^s)
=-(2\kappa-1)\sum_{j=1}^{\kappa}\frac{1}{1+e^s\delta_j}\;<\;0\;}
\]
从而存在有限极限
\[
\boxed{\;
\lim_{m\to\infty}\Xi(m)
=\Xi(1)-(2\kappa-1)\sum_{j=1}^{\kappa}\log\frac{1+\delta_j}{\delta_j}\; }.
\]
\end{corollary}
\begin{proof}
由命题 \ref{prop:xi-hankel-det-scaling-covariance}，
\(
S_{\mathrm{bulk}}(m)=S_{\mathrm{bulk}}(1)-\log J_{\kappa-1}(m;\delta)
\)；
而由命题 \ref{prop:xi-pick-poisson-det-scaling-law}，
\(
S_{\mathrm{gen}}(m)=S_{\mathrm{gen}}(1)+\kappa\log m
\)。
因此
\[
\Xi(m)
=\Xi(1)-\log J_{\kappa-1}(m;\delta)-\kappa^2\log m.
\]
将
\(
\log J_{\kappa-1}(m;\delta)
=\kappa(\kappa-1)\log m+(2\kappa-1)\sum_{j=1}^{\kappa}\bigl(\log(1+\delta_j)-\log(1+m\delta_j)\bigr)
\)
代入并整理得
\[
\Xi(m)=\Xi(1)+(2\kappa-1)\sum_{j=1}^{\kappa}\log\frac{1+m\delta_j}{m(1+\delta_j)}.
\]
令 \(m=e^s\)，对 \(s\) 求导即得
\[
\frac{\dd}{\dd s}\,\Xi(e^s)
=(2\kappa-1)\sum_{j=1}^{\kappa}\left(\frac{e^s\delta_j}{1+e^s\delta_j}-1\right)
=-(2\kappa-1)\sum_{j=1}^{\kappa}\frac{1}{1+e^s\delta_j}\;<0.
\]
最后由 \(\log(1+m\delta_j)=\log m+\log\delta_j+o(1)\)（\(m\to\infty\)）得到极限式。证毕。
\end{proof}

\begin{corollary}[bulk 尺度熵流的闭式：logistic 阶梯分解与 \texorpdfstring{$\varphi$}{phi}-归一有效维度]\label{cor:xi-bulk-entropy-rg-staircase}
\leanverified{paper\_xi\_bulk\_entropy\_rg\_staircase}
仍在命题 \ref{prop:xi-hankel-det-scaling-covariance} 的有限原子设定下，令 \(s=\log m\)。
则 \(S_{\mathrm{bulk}}(e^s)\) 在 \(s\in\RR\) 上可微，并满足一阶流方程
\[
\boxed{\;
\frac{\dd}{\dd s}S_{\mathrm{bulk}}(e^s)
=\kappa^{2}-(2\kappa-1)\sum_{j=1}^{\kappa}\frac{1}{1+e^{s}\delta_{j}}\;}
\]
从而其 \(\varphi\)-标度归一的尺度有效维度
\[
d_{\mathrm{eff}}^{\mathrm{bulk}}(s):=\frac{1}{\log\varphi}\frac{\dd}{\dd s}S_{\mathrm{bulk}}(e^s)
\]
是 \(\kappa\) 个平移 logistic 尾函数 \((1+e^{s}\delta_j)^{-1}\) 的线性叠加。
每个深度原子 \(\delta_j\) 诱导一处以 \(\sigma_j:=-\log\delta_j\) 为中心的单调跃迁台阶，其台阶高度恒为 \((2\kappa-1)/\log\varphi\)，而端点值仅由 \(\kappa\) 决定。
\end{corollary}
\begin{proof}
由 \(S_{\mathrm{bulk}}(m)=S_{\mathrm{bulk}}(1)-\log J_{\kappa-1}(m;\delta)\) 与命题 \ref{prop:xi-hankel-det-scaling-covariance} 的闭式
\[
\log J_{\kappa-1}(m;\delta)
=\kappa(\kappa-1)\log m
+(2\kappa-1)\sum_{j=1}^{\kappa}\bigl(\log(1+\delta_j)-\log(1+m\delta_j)\bigr),
\]
令 \(m=e^s\) 并对 \(s\) 求导，得到
\[
\frac{\dd}{\dd s}S_{\mathrm{bulk}}(e^s)
=-\kappa(\kappa-1)+(2\kappa-1)\sum_{j=1}^{\kappa}\frac{e^s\delta_j}{1+e^s\delta_j}
=\kappa^2-(2\kappa-1)\sum_{j=1}^{\kappa}\frac{1}{1+e^s\delta_j}.
\]
其余表述由 \(\frac{1}{1+e^s\delta_j}=\frac{1}{1+e^{s-\sigma_j}}\) 与线性叠加直接得到。证毕。
\end{proof}

\begin{corollary}[bulk 熵的严格凸性：伸缩流的熵曲率闭式且处处正]\label{cor:xi-bulk-entropy-strict-convexity}
\leanverified{paper\_xi\_bulk\_entropy\_strict\_convexity}
仍在命题 \ref{prop:xi-hankel-det-scaling-covariance} 的有限原子设定下，令 \(s=\log m\)。
则 \(S_{\mathrm{bulk}}(e^s)\) 二阶可导并满足严格凸性
\[
\boxed{\;
\frac{\dd^2}{\dd s^2}S_{\mathrm{bulk}}(e^s)
=(2\kappa-1)\sum_{j=1}^{\kappa}\frac{e^s\delta_j}{(1+e^s\delta_j)^2}\;>\;0\;}
\qquad(s\in\RR).
\]
\end{corollary}
\begin{proof}
由 \(S_{\mathrm{bulk}}(m)=S_{\mathrm{bulk}}(1)-\log J_{\kappa-1}(m;\delta)\) 与命题 \ref{prop:xi-hankel-det-scaling-covariance} 的闭式，
\[
S_{\mathrm{bulk}}(e^s)
=S_{\mathrm{bulk}}(1)-\kappa(\kappa-1)s-(2\kappa-1)\sum_{j=1}^{\kappa}\bigl(\log(1+\delta_j)-\log(1+e^s\delta_j)\bigr).
\]
两次求导并化简即得所示曲率公式；各项严格为正，故严格凸。证毕。
\end{proof}

% (User input window)

\begin{remark}[bulk 熵势的费米化解释：两态模的巨正则形式与半填充峰]\label{rem:xi-bulk-entropy-fermi-grand-canonical}
沿用命题 \ref{prop:xi-hankel-det-scaling-covariance} 的有限原子记号并令 \(m=e^s\)。
定义单模配分因子与占据函数
\[
Z_j(s):=1+e^{s}\delta_j,
\qquad
n_j(s):=\frac{e^{s}\delta_j}{1+e^{s}\delta_j}\in(0,1).
\]
则由推论 \ref{cor:xi-bulk-entropy-strict-convexity} 的证明中闭式可将 bulk 熵势写成
\[
\boxed{\;
S_{\mathrm{bulk}}(e^{s})
=S_{\mathrm{bulk}}(1)-\kappa(\kappa-1)s
+(2\kappa-1)\sum_{j=1}^{\kappa}\log\!\Bigl(\frac{Z_j(s)}{Z_j(0)}\Bigr)\;}
\]
亦即
\(
S_{\mathrm{bulk}}(e^{s})=\mathrm{const}-\kappa(\kappa-1)s+(2\kappa-1)\sum_j\log Z_j(s)
\)。
因此 \(S_{\mathrm{bulk}}\) 在严格意义上等同于 \(\kappa\) 个独立两态模（费米占据）巨正则自由能的线性重标度：\(\log\delta_j\) 作为一体偏置进入，\(s\) 作为线性控制参量进入（可按物理习惯解释为化学势坐标或等价位势坐标）。
\par
由推论 \ref{cor:xi-bulk-entropy-rg-staircase} 与推论 \ref{cor:xi-bulk-entropy-strict-convexity}，一阶与二阶导数可写为
\[
\frac{\dd}{\dd s}S_{\mathrm{bulk}}(e^s)=-\kappa(\kappa-1)+(2\kappa-1)\sum_{j=1}^{\kappa}n_j(s),
\qquad
\frac{\dd^2}{\dd s^2}S_{\mathrm{bulk}}(e^s)=(2\kappa-1)\sum_{j=1}^{\kappa}n_j(s)\bigl(1-n_j(s)\bigr).
\]
其中每个涨落项 \(n_j(1-n_j)\) 在 \(s=\sigma_j:=-\log\delta_j\) 处达到唯一峰值，并满足 \(n_j(\sigma_j)=1/2\)。
\end{remark}

\begin{proposition}[bulk 费米化势的微正则对偶：Legendre 变换给出二元熵和表示且无连续相共存]\label{prop:xi-bulk-entropy-microcanonical-legendre-dual}
在注 \ref{rem:xi-bulk-entropy-fermi-grand-canonical} 的记号下，令
\[
\Phi_{\mathrm{F}}(s):=\sum_{j=1}^{\kappa}\log(1+e^s\delta_j),
\qquad
N(s):=\Phi_{\mathrm{F}}'(s)=\sum_{j=1}^{\kappa}n_j(s).
\]
则 \(N(s)\) 严格递增并将 \(s\in\RR\) 双射到 \((0,\kappa)\)。
对任意 \(N\in(0,\kappa)\)，定义 Legendre 对偶（微正则熵）
\[
\Sigma(N):=\sup_{s\in\RR}\bigl(sN-\Phi_{\mathrm{F}}(s)\bigr).
\]
则最优点 \(s(N)\) 唯一。
令 \(n_j:=n_j(s(N))\)，则存在完全显式的二元熵和表示
\[
\boxed{\;
\Sigma(N)= -\sum_{j=1}^{\kappa}\Bigl(n_j\log n_j+(1-n_j)\log(1-n_j)\Bigr)\;-\;\sum_{j=1}^{\kappa}n_j\log\delta_j\;.}
\]
并且由严格凸性，\(\Sigma\) 不存在平坦段，从而以 \(s\) 为连续控制参量时不发生连续相共存。
\end{proposition}
\begin{proof}
由
\(
\Phi_{\mathrm{F}}''(s)=\sum_{j=1}^{\kappa}\frac{e^s\delta_j}{(1+e^s\delta_j)^2}
=\sum_{j=1}^{\kappa}n_j(s)\bigl(1-n_j(s)\bigr)>0
\)
知 \(\Phi_{\mathrm{F}}\) 严格凸，故 \(N=\Phi_{\mathrm{F}}'\) 严格递增。
并且当 \(s\to-\infty\) 时 \(n_j(s)\to 0\)、当 \(s\to+\infty\) 时 \(n_j(s)\to 1\)，从而 \(N(s)\) 将 \(\RR\) 双射到 \((0,\kappa)\)。
因此对每个 \(N\in(0,\kappa)\) 存在唯一 \(s(N)\) 使 \(N=N(s(N))\)，并给出 Legendre 极值点。
\par
令 \(s=s(N)\) 且 \(n_j=n_j(s)\)。
由 \(n_j/(1-n_j)=e^s\delta_j\) 得
\(
s=\log\frac{n_j}{1-n_j}-\log\delta_j
\)，
并且
\(
\log(1+e^s\delta_j)=-\log(1-n_j)
\)。
故
\[
sN-\Phi_{\mathrm{F}}(s)
=\sum_{j=1}^{\kappa}\Bigl(s\,n_j-\log(1+e^s\delta_j)\Bigr)
=-\sum_{j=1}^{\kappa}\Bigl(n_j\log n_j+(1-n_j)\log(1-n_j)\Bigr)-\sum_{j=1}^{\kappa}n_j\log\delta_j,
\]
得到所示闭式。
最后，严格凸性排除 \(\Phi_{\mathrm{F}}\) 的线性段，从而 \(\Sigma\) 不可能出现对应的平坦段；因此不存在由 \(s\)-连续奇性诱导的相共存。证毕。
\end{proof}

% (User input window)

推论 \ref{cor:xi-bulk-entropy-strict-convexity} 给出的二阶尺度响应具有固定核形状，从而 bulk 熵的曲率在尺度轴上形成一条可层析、可归一并可识别的深度谱接口。

\begin{proposition}[尺度熵曲率的 logistic 核分解：临界尺度 \texorpdfstring{$\sigma_j=-\log\delta_j$}{sigma_j=-log delta_j}]\label{prop:xi-bulk-curvature-logistic-kernel-decomposition}
\leanverified{paper\_xi\_bulk\_curvature\_logistic\_kernel\_decomposition}
在推论 \ref{cor:xi-bulk-entropy-strict-convexity} 的有限原子设定下，定义 bulk 尺度熵曲率密度
\[
\mathcal{K}_{\mathrm{bulk}}(s):=\frac{\dd^{2}}{\dd s^{2}}S_{\mathrm{bulk}}(e^{s})\qquad(s\in\RR),
\]
并引入标准 logistic 密度
\[
\phi(t):=\frac{e^{t}}{(1+e^{t})^{2}}\qquad(t\in\RR),
\]
以及临界尺度 \(\sigma_j:=-\log\delta_j\)。
则有严格分解
\[
\boxed{\;
\frac{1}{2\kappa-1}\mathcal{K}_{\mathrm{bulk}}(s)=\sum_{j=1}^{\kappa}\phi(s-\sigma_j).\;}
\]
并且 \(\phi\) 在 \(\RR\) 上单峰：其唯一最大值为 \(\phi(0)=\tfrac14\)，且对任意 \(t\neq 0\) 有 \(\phi(t)<\tfrac14\)。
因此 \(\{\sigma_j\}\) 给出曲率核族在尺度轴上的平移中心，对应统一伸缩 \(m=e^s\) 下的临界尺度 \(m\asymp \delta_j^{-1}\)。
\end{proposition}
\begin{proof}
由推论 \ref{cor:xi-bulk-entropy-strict-convexity}，
\[
\mathcal{K}_{\mathrm{bulk}}(s)
=(2\kappa-1)\sum_{j=1}^{\kappa}\frac{e^{s}\delta_j}{(1+e^{s}\delta_j)^{2}}.
\]
记 \(\delta_j=e^{-\sigma_j}\)，则 \(\frac{e^{s}\delta_j}{(1+e^{s}\delta_j)^{2}}=\phi(s-\sigma_j)\)，得到所示分解。
另一方面，
\[
\phi'(t)=\frac{e^{t}(1-e^{t})}{(1+e^{t})^{3}},
\]
故 \(\phi'(t)>0\) 当 \(t<0\)、\(\phi'(t)<0\) 当 \(t>0\)，从而 \(\phi\) 单峰且最大值在 \(t=0\) 唯一取得；\(\phi(0)=\tfrac14\) 由代入得到。证毕。
\end{proof}

\begin{corollary}[熵曲率的总质量量子化：\texorpdfstring{$\int\mathcal{K}_{\mathrm{bulk}}=\kappa(2\kappa-1)$}{int K_bulk = k(2k-1)}]\label{cor:xi-bulk-curvature-mass-quantization}
在命题 \ref{prop:xi-bulk-curvature-logistic-kernel-decomposition} 的记号下，有
\[
\boxed{\;
\int_{-\infty}^{\infty}\mathcal{K}_{\mathrm{bulk}}(s)\,\dd s=\kappa(2\kappa-1),
\qquad
\int_{-\infty}^{\infty}\phi(s-\sigma_j)\,\dd s=1\ \ (1\le j\le\kappa).\;}
\]
\end{corollary}
\begin{proof}
对任意固定 \(j\)，令 \(u=e^{s-\sigma_j}\)，则 \(\dd s=\dd u/u\)，并且
\[
\int_{-\infty}^{\infty}\phi(s-\sigma_j)\,\dd s
=\int_{0}^{\infty}\frac{u}{(1+u)^2}\frac{\dd u}{u}
=\int_{0}^{\infty}\frac{\dd u}{(1+u)^2}
=1.
\]
将其代入命题 \ref{prop:xi-bulk-curvature-logistic-kernel-decomposition} 的分解式即得曲率总质量公式。证毕。
\end{proof}

\begin{proposition}[熵曲率层析的可识别性：由 \texorpdfstring{$\mathcal{K}_{\mathrm{bulk}}$}{K_bulk} 唯一重建深度谱]\label{prop:xi-bulk-curvature-tomography-identifiability}
\leanverified{paper\_xi\_bulk\_curvature\_tomography\_identifiability}
\leanverified{paper\_xi\_bulk\_curvature\_tomography\_identifiability}
固定 \(\kappa\) 并设 \(\mathcal{K}_{\mathrm{bulk}}:\RR\to\RR\) 已知为连续函数，且满足命题 \ref{prop:xi-bulk-curvature-logistic-kernel-decomposition} 的有限原子分解。
则多重集 \(\{\sigma_j\}_{j=1}^{\kappa}\)（从而 \(\{\delta_j=e^{-\sigma_j}\}_{j=1}^{\kappa}\)）由 \(\mathcal{K}_{\mathrm{bulk}}\) 唯一确定（忽略排列）。
\end{proposition}
\begin{proof}
取 Fourier 变换
\[
\widehat f(\xi):=\int_{\RR}e^{-\mathrm{i}\xi s}f(s)\,\dd s.
\]
记 \(\widehat\phi(0)=1\) 为连续延拓值，则 \(\phi\) 的 Fourier 变换具有经典闭式
\[
\widehat\phi(\xi)=\frac{\pi\xi}{\sinh(\pi\xi)}\qquad(\xi\in\RR),
\qquad
\widehat\phi(\xi)\neq 0\quad(\xi\in\RR).
\]
对命题 \ref{prop:xi-bulk-curvature-logistic-kernel-decomposition} 两边取 Fourier 变换得
\[
\sum_{j=1}^{\kappa}e^{-\mathrm{i}\xi\sigma_j}
=\frac{1}{\widehat\phi(\xi)}\cdot
\widehat{\Bigl(\frac{1}{2\kappa-1}\mathcal{K}_{\mathrm{bulk}}\Bigr)}(\xi)
\qquad(\xi\in\RR).
\]
右端由 \(\mathcal{K}_{\mathrm{bulk}}\) 决定，故得到离散测度 \(\nu:=\sum_{j=1}^{\kappa}\delta_{\sigma_j}\) 的 Fourier 变换 \(\widehat\nu(\xi)=\sum_j e^{-\mathrm{i}\xi\sigma_j}\)。
Fourier 变换在有限 Radon 测度上为单射，因此 \(\nu\)（等价多重集 \(\{\sigma_j\}\)）唯一确定；从而 \(\delta_j=e^{-\sigma_j}\) 唯一确定。证毕。
\end{proof}

\begin{remark}[全息差额的闭式尺度轨迹：\texorpdfstring{$\Xi(e^s)$}{Xi(e^s)} 的等价表达]\label{rem:xi-holographic-gap-closed-trajectory}
推论 \ref{cor:xi-bulk-boundary-holographic-gap-dissipation} 的闭式可在对数尺度变量 \(s=\log m\) 下改写为
\[
\Xi(e^s)
=\Xi(1)-(2\kappa-1)\sum_{j=1}^{\kappa}\log\!\Bigl(\frac{e^{s}(1+\delta_j)}{1+e^{s}\delta_j}\Bigr),
\]
从而 \(\Xi(e^s)\) 的全部尺度依赖由 \(\{\delta_j\}\) 的 logistic 尾项封闭。
\end{remark}

\begin{proposition}[伸缩流的速度匹配：bulk 与 boundary 的熵增长率在唯一尺度相交]\label{prop:xi-bulk-boundary-speed-matching}
\leanverified{paper\_xi\_bulk\_boundary\_speed\_matching}
设 \(\kappa\ge 2\)。
定义
\[
v_{\mathrm{bulk}}(s):=\frac{\dd}{\dd s}S_{\mathrm{bulk}}(e^s),
\qquad
v_{\mathrm{gen}}(s):=\frac{\dd}{\dd s}S_{\mathrm{gen}}(e^s)=\kappa
\]
（其中等号由推论 \ref{cor:xi-pick-poisson-entropy-derivative-visible-dimension} 给出）。
则存在唯一 \(s_\star\in\RR\) 使 \(v_{\mathrm{bulk}}(s_\star)=v_{\mathrm{gen}}(s_\star)\)。
该唯一尺度满足显式方程
\[
\boxed{\;
\sum_{j=1}^{\kappa}\frac{1}{1+e^{s_\star}\delta_j}=\frac{\kappa(\kappa-1)}{2\kappa-1}.\;}
\]
\end{proposition}
\begin{proof}
由推论 \ref{cor:xi-bulk-entropy-strict-convexity}，\(v_{\mathrm{bulk}}'(s)=\mathcal{K}_{\mathrm{bulk}}(s)>0\)，故 \(v_{\mathrm{bulk}}\) 在 \(\RR\) 上严格递增。
由推论 \ref{cor:xi-bulk-entropy-rg-staircase} 的闭式，
\[
\lim_{s\to-\infty}v_{\mathrm{bulk}}(s)=\kappa^2-(2\kappa-1)\kappa=-\kappa(\kappa-1),
\qquad
\lim_{s\to\infty}v_{\mathrm{bulk}}(s)=\kappa^2.
\]
当 \(\kappa\ge 2\) 时有 \(-\kappa(\kappa-1)<\kappa<\kappa^2\)，故由介值定理存在 \(s_\star\) 使 \(v_{\mathrm{bulk}}(s_\star)=\kappa\)；严格递增性给出唯一性。
将 \(v_{\mathrm{bulk}}(s_\star)=\kappa\) 代入推论 \ref{cor:xi-bulk-entropy-rg-staircase} 的闭式并整理即得所示方程。证毕。
\end{proof}

\begin{corollary}[速度匹配尺度的均值夹逼：由 \texorpdfstring{$\delta_{\min}$}{delta_min} 与 \texorpdfstring{$\bar\delta$}{delta_bar} 给出的无参数区间]\label{cor:xi-bulk-boundary-speed-matching-scale-bracketing}
令 \(s_\star\) 为命题 \ref{prop:xi-bulk-boundary-speed-matching} 的唯一解，并记
\[
\delta_{\min}:=\min_{1\le j\le\kappa}\delta_j,
\qquad
\bar\delta:=\frac{1}{\kappa}\sum_{j=1}^{\kappa}\delta_j.
\]
则 \(e^{s_\star}\) 满足显式夹逼
\[
\boxed{\;
\frac{\kappa}{(\kappa-1)\,\bar\delta}\ \le\ e^{s_\star}\ \le\ \frac{\kappa}{(\kappa-1)\,\delta_{\min}}\;}
\]
从而
\[
\log\frac{\kappa}{\kappa-1}-\log\bar\delta\ \le\ s_\star\ \le\ \log\frac{\kappa}{\kappa-1}-\log\delta_{\min}.
\]
因此即便不恢复全体 \(\{\delta_j\}\)，仅凭 \((\bar\delta,\delta_{\min})\) 亦可把 bulk--boundary 的速度交汇尺度压缩到一个完全可审计的无参数区间。
\end{corollary}
\begin{proof}
将命题 \ref{prop:xi-bulk-boundary-speed-matching} 的方程除以 \(\kappa\) 得
\[
\frac{1}{\kappa}\sum_{j=1}^{\kappa}\frac{1}{1+e^{s_\star}\delta_j}=\frac{\kappa-1}{2\kappa-1}.
\]
对固定 \(a>0\)，函数 \(x\mapsto (1+ax)^{-1}\) 在 \((0,\infty)\) 上单调递减且凸。
故由 Jensen 不等式得
\[
\frac{1}{1+a\bar\delta}\le \frac{1}{\kappa}\sum_{j=1}^{\kappa}\frac{1}{1+a\delta_j},
\]
令 \(a=e^{s_\star}\) 并代入右端的显式值，整理得
\(
e^{s_\star}\ge \frac{\kappa}{(\kappa-1)\bar\delta}
\)。
另一方面，因 \(\delta_j\ge \delta_{\min}\) 且函数递减，
\[
\frac{1}{\kappa}\sum_{j=1}^{\kappa}\frac{1}{1+a\delta_j}\le \frac{1}{1+a\delta_{\min}}.
\]
同样取 \(a=e^{s_\star}\) 并代入，整理得
\(
e^{s_\star}\le \frac{\kappa}{(\kappa-1)\delta_{\min}}
\)。
对数形式由两端取对数得到。证毕。
\end{proof}

\begin{corollary}[曲率概率化：归一曲率给出尺度分布，其矩分解为普适项与深度谱项]\label{cor:xi-bulk-curvature-probability-moments}
在命题 \ref{prop:xi-bulk-curvature-logistic-kernel-decomposition} 的记号下，令
\[
p(s):=\frac{1}{\kappa(2\kappa-1)}\mathcal{K}_{\mathrm{bulk}}(s).
\]
则 \(p\) 为 \(\RR\) 上概率密度，并具有等权混合表示
\[
\boxed{\;
p(s)=\frac{1}{\kappa}\sum_{j=1}^{\kappa}\phi(s-\sigma_j).\;}
\]
记 \(\bar\sigma:=\frac{1}{\kappa}\sum_{j=1}^\kappa\sigma_j\) 并定义离散方差
\[
\Var(\sigma):=\frac{1}{\kappa}\sum_{j=1}^{\kappa}(\sigma_j-\bar\sigma)^2.
\]
则一阶矩与二阶中心矩满足分解
\[
\boxed{\;
\int_{\RR}s\,p(s)\,\dd s=\bar\sigma,\qquad
\int_{\RR}(s-\bar\sigma)^2\,p(s)\,\dd s=\frac{\pi^2}{3}+\Var(\sigma).\;}
\]
\end{corollary}
\begin{proof}
由推论 \ref{cor:xi-bulk-curvature-mass-quantization} 得 \(\int_{\RR}p(s)\dd s=1\)，故 \(p\) 为概率密度；混合表示由命题 \ref{prop:xi-bulk-curvature-logistic-kernel-decomposition} 直接得到。
由于 \(\phi\) 为中心 logistic 分布密度，满足 \(\int_{\RR} t\,\phi(t)\dd t=0\) 与 \(\int_{\RR} t^2\,\phi(t)\dd t=\pi^2/3\)。
对混合分布把变量平移 \(s=t+\sigma_j\) 后逐项积分即可得到均值与方差分解。证毕。
\end{proof}

\begin{proposition}[熵曲率分布的矩母函数因式分解：普适 \texorpdfstring{$\zeta(2n)$}{zeta(2n)} 塔与深度谱累积量分裂]\label{prop:xi-bulk-curvature-mgf-factorization-zeta-tower}
令 \(p\) 与 \(\{\sigma_j\}\) 如推论 \ref{cor:xi-bulk-curvature-probability-moments}。
则存在解析邻域 \(|t|<1\) 使矩母函数严格因式分解为
\[
\boxed{\;
\int_{\RR}e^{ts}\,p(s)\,\dd s
=\Gamma(1+t)\Gamma(1-t)\cdot \frac{1}{\kappa}\sum_{j=1}^{\kappa}e^{t\sigma_j}\;}
\qquad(|t|<1).
\]
由 Euler 反射与 \(\sin\) 的 Hadamard 乘积，
\[
\Gamma(1+t)\Gamma(1-t)=\frac{\pi t}{\sin(\pi t)}
=\prod_{n\ge 1}\Bigl(1-\frac{t^2}{n^2}\Bigr)^{-1},
\]
从而累积量母函数在 \(|t|<1\) 上严格分裂为
\[
\log\!\int_{\RR}e^{ts}\,p(s)\,\dd s
=\sum_{n\ge 1}\frac{\zeta(2n)}{n}\,t^{2n}
\;+\;\log\!\Bigl(\frac{1}{\kappa}\sum_{j=1}^{\kappa}e^{t\sigma_j}\Bigr).
\]
记 \(\kappa_q(p):=\left.\frac{\dd^q}{\dd t^q}\log\int_{\RR}e^{ts}p(s)\dd s\right|_{t=0}\) 为 \(p\) 的 \(q\) 阶累积量，
并记 \(\kappa_q(\sigma):=\left.\frac{\dd^q}{\dd t^q}\log\bigl(\kappa^{-1}\sum_j e^{t\sigma_j}\bigr)\right|_{t=0}\) 为等权离散测度 \(\kappa^{-1}\sum_j\delta_{\sigma_j}\) 的累积量。
则对一切 \(n\ge 1\) 成立
\[
\boxed{\;
\kappa_{2n}(p)=(2n)!\,\frac{\zeta(2n)}{n}\;+\;\kappa_{2n}(\sigma),
\qquad
\kappa_{2n+1}(p)=\kappa_{2n+1}(\sigma).\;}
\]
因此 bulk 尺度响应的高阶常数被严格锁定为偶 \(\zeta\) 值与深度谱离散累积量两部分之和。
\end{proposition}
\begin{proof}
由推论 \ref{cor:xi-bulk-curvature-probability-moments}，
\(
p(s)=\kappa^{-1}\sum_{j=1}^\kappa \phi(s-\sigma_j)
\)。
故在 \(|t|<1\) 内可交换求和与积分并得
\[
\int_{\RR}e^{ts}\,p(s)\,\dd s
=\frac{1}{\kappa}\sum_{j=1}^{\kappa}e^{t\sigma_j}\int_{\RR}e^{tu}\phi(u)\,\dd u.
\]
对 \(\phi(u)=\frac{e^{u}}{(1+e^{u})^{2}}\) 作代换 \(y=e^{u}\)（\(\dd u=\dd y/y\)），得到
\[
\int_{\RR}e^{tu}\phi(u)\,\dd u
=\int_{0}^{\infty}\frac{y^{t}}{(1+y)^{2}}\,\dd y
=\mathrm{B}(1+t,1-t)
=\Gamma(1+t)\Gamma(1-t),
\]
其中 Beta 恒等式要求 \(-1<\Re(t)<1\)，故在实轴给出 \(|t|<1\) 的解析域。
从而得第一条因式分解。
第二条分裂由 Euler 反射公式与 \(\sin\) 的乘积展开（或等价地由 \(\log\Gamma(1\pm t)\) 的 \(\zeta\)-级数）得到。
最后对两边在 \(t=0\) 处求导，利用
\(
\frac{\dd^{2n}}{\dd t^{2n}}\sum_{m\ge 1}\frac{\zeta(2m)}{m}t^{2m}\big|_{t=0}
=(2n)!\,\frac{\zeta(2n)}{n}
\)
并注意奇次项消失，得到所示累积量分裂。证毕。
\end{proof}

\begin{proposition}[\texorpdfstring{$\Gamma$}{Gamma}-剥离的深度谱显影：固定解析乘子将 bulk 曲率降维为有限指数和]\label{prop:xi-gamma-stripped-depth-spectrum-imaging}
在命题 \ref{prop:xi-bulk-curvature-mgf-factorization-zeta-tower} 的记号下，定义
\[
\mathcal M_{\mathrm{depth}}(t):=\frac{\sin(\pi t)}{\pi t}\int_{\RR}e^{ts}\,p(s)\,\dd s\qquad(|t|<1).
\]
则在 \(|t|<1\) 上严格有
\[
\boxed{\;
\mathcal M_{\mathrm{depth}}(t)=\frac{1}{\kappa}\sum_{j=1}^{\kappa}e^{t\sigma_j}.\;}
\]
右端为指数多项式，因而 \(\mathcal M_{\mathrm{depth}}\) 作为整函数延拓至全平面，且其增长型由 \(\max_j|\sigma_j|\) 唯一封口。
并且：
\begin{enumerate}
  \item \(\mathcal M_{\mathrm{depth}}\) 在任意非空开区间上的取值即唯一确定多重集 \(\{\sigma_j\}\)（指数多项式的解析唯一性）；
  \item \(\mathcal M_{\mathrm{depth}}^{(n)}(0)=\kappa^{-1}\sum_{j}\sigma_j^{n}\) 给出全部幂和，结合 Newton 恒等式可在有限步内恢复特征多项式 \(\prod_{j=1}^{\kappa}(z-\sigma_j)\)；
  \item 由于 \(\sin(\pi t)/(\pi t)\) 在 \(|t|<1\) 内无零点，\(\int e^{ts}p(s)\dd s\) 的可见零点结构在剥离后完全由有限指数和 \(\sum_j e^{t\sigma_j}\) 决定而不受普适核形状干扰。
\end{enumerate}
该结论给出一条可审计的“显影接口”：bulk 的尺度曲率经一个固定的解析乘子即可无损还原 finite-rank 深度谱。
\end{proposition}
\begin{proof}
由命题 \ref{prop:xi-bulk-curvature-mgf-factorization-zeta-tower}，
\(
\int e^{ts}p(s)\dd s=\frac{\pi t}{\sin(\pi t)}\cdot \kappa^{-1}\sum_j e^{t\sigma_j}
\)
在 \(|t|<1\) 成立；两边同乘 \(\sin(\pi t)/(\pi t)\) 即得所示恒等式。
指数多项式的整性与增长型为经典事实。
条目 (1)--(2) 分别由解析延拓唯一性与幂和--系数（Newton）恒等式给出；
条目 (3) 由剥离因子在 \(|t|<1\) 内无零点与恒等式右端决定性得到。证毕。
\end{proof}

\begin{corollary}[曲率上界与 Lipschitz 速度极限：bulk 尺度响应的统一封口]\label{cor:xi-bulk-curvature-max-lipschitz}
在命题 \ref{prop:xi-bulk-curvature-logistic-kernel-decomposition} 的记号下，有点态上界
\[
\boxed{\;
0<\mathcal{K}_{\mathrm{bulk}}(s)\le \frac{(2\kappa-1)\kappa}{4}\qquad(s\in\RR).\;}
\]
从而 \(v_{\mathrm{bulk}}(s)=\frac{\dd}{\dd s}S_{\mathrm{bulk}}(e^s)\) 在 \(\RR\) 上满足统一 Lipschitz 约束
\[
\boxed{\;
|v_{\mathrm{bulk}}(s_2)-v_{\mathrm{bulk}}(s_1)|
\le \frac{(2\kappa-1)\kappa}{4}\,|s_2-s_1|\qquad(s_1,s_2\in\RR).\;}
\]
该上界与 \(\{\delta_j\}\) 无关，仅由 \(\kappa\) 与系数 \((2\kappa-1)\) 决定。
\end{corollary}
\begin{proof}
由命题 \ref{prop:xi-bulk-curvature-logistic-kernel-decomposition}，
\(
\mathcal{K}_{\mathrm{bulk}}(s)=(2\kappa-1)\sum_{j=1}^{\kappa}\phi(s-\sigma_j)
\)；
而 \(\phi(t)\le \tfrac14\)，故得点态上界。
由于 \(v_{\mathrm{bulk}}'(s)=\mathcal{K}_{\mathrm{bulk}}(s)\) 且 \(\mathcal{K}_{\mathrm{bulk}}\) 有界，积分估计即得 Lipschitz 约束。证毕。
\end{proof}

\begin{remark}[尺度峰的可分辨性：分离深度谱给出曲率层析的稳定分解]\label{rem:xi-bulk-curvature-peak-resolution}
对 \(\phi\) 记其分布函数 \(F(t)=\frac{1}{1+e^{-t}}\)。则对任意 \(T>0\),
\[
\int_{|t|\le T}\phi(t)\,\dd t
=F(T)-F(-T)
=1-\frac{2}{1+e^{T}}.
\]
因此在固定容差 \(\varepsilon\in(0,1)\) 下，取 \(T_\varepsilon:=\log\!\bigl(\frac{2}{\varepsilon}-1\bigr)\) 有
\(
\int_{|t|>T_\varepsilon}\phi(t)\dd t\le \varepsilon
\)。
若最小尺度分离度满足 \(\min_{j\neq k}|\sigma_j-\sigma_k|\ge 2T_\varepsilon\)，则区间 \([\sigma_j-T_\varepsilon,\sigma_j+T_\varepsilon]\) 两两不交，
且每个区间承载 \(\phi(s-\sigma_j)\) 的至少 \(1-\varepsilon\) 质量，从而 \(\mathcal{K}_{\mathrm{bulk}}\) 在制度精度 \(\varepsilon\) 下分解为 \(\kappa\) 个互不混叠的局部峰包络。
反之，若存在一族 \(\sigma_j\) 在尺度窗内簇化，则对应核在该窗内发生线性叠加并抬升峰高（重数 \(r\) 的极限峰值为 \(r(2\kappa-1)/4\)），体现为可观测的峰合并。
\end{remark}

除上述以峰分离为核心的层析口径外，\(\phi\) 的平移族在 \(L^2(\RR)\) 中还诱导一组可计算的相干几何量；其局部体积结构将把偶 \(\zeta\) 常数以确定方式嵌入到尺度几何中。

\begin{proposition}[尺度相干核的闭式：\texorpdfstring{$\mathrm{sech}^2$}{sech^2} 平移族的 Gram 体积与双曲函数显式化]\label{prop:xi-scale-coherence-kernel-closed-form}
\leanverified{paper\_xi\_scale\_coherence\_kernel\_closed\_form}
将 logistic 核写为等价形式
\[
\phi(t)=\frac{1}{4\cosh^2(t/2)}=\frac{1}{4}\,\mathrm{sech}^2(t/2).
\]
定义两点尺度相干核（平移族的 \(L^2\) 内积）
\[
K(\Delta):=\int_{\RR}\phi(s)\phi(s-\Delta)\,\dd s\qquad(\Delta\in\RR).
\]
则 \(K\) 为偶函数，并存在严格闭式
\[
\boxed{\;
K(\Delta)=\frac{\Delta\cosh(\Delta/2)-2\sinh(\Delta/2)}{4\sinh^3(\Delta/2)}\quad(\Delta\neq 0),
\qquad
K(0)=\frac{1}{6}.\;}
\]
其在 \(\Delta=0\) 的局部展开为
\[
K(\Delta)=\frac{1}{6}-\frac{\Delta^2}{60}+\frac{\Delta^4}{1008}+O(\Delta^6).
\]
\end{proposition}
\begin{proof}
令 \(a:=\Delta/2\) 并作变量替换 \(s=2x\)。由 \(\phi(2x)=\frac14\mathrm{sech}^2(x)\) 得
\[
K(\Delta)=\frac{1}{16}\int_{\RR}\mathrm{sech}^2(x)\,\mathrm{sech}^2(x-a)\,2\,\dd x
=\frac18\int_{\RR}\mathrm{sech}^2(x)\,\mathrm{sech}^2(x-a)\,\dd x.
\]
令 \(y=\tanh x\)，则 \(\dd y=\mathrm{sech}^2(x)\dd x\)，并记 \(t:=\tanh a\)。
利用双曲正切差公式
\(
\tanh(x-a)=\frac{y-t}{1-ty}
\)
得到
\[
\mathrm{sech}^2(x-a)=1-\tanh^2(x-a)=\mathrm{sech}^2(a)\cdot\frac{1-y^2}{(1-ty)^2}.
\]
从而
\[
\int_{\RR}\mathrm{sech}^2(x)\,\mathrm{sech}^2(x-a)\,\dd x
=\mathrm{sech}^2(a)\int_{-1}^{1}\frac{1-y^2}{(1-ty)^2}\,\dd y.
\]
直接积分（或令 \(u=1-ty\)）得
\[
\int_{-1}^{1}\frac{1-y^2}{(1-ty)^2}\,\dd y
=\frac{2}{t^3}\Bigl(\log\frac{1+t}{1-t}-2t\Bigr)
=\frac{4}{t^3}(a-t),
\]
其中使用 \(\log\frac{1+t}{1-t}=2\operatorname{arctanh}(t)=2a\)。
于是
\[
K(\Delta)
=\frac18\cdot \mathrm{sech}^2(a)\cdot \frac{4}{t^3}(a-t)
=\frac12\cdot\frac{a\cosh a-\sinh a}{\sinh^3 a},
\]
代回 \(a=\Delta/2\) 即得 \(\Delta\neq 0\) 的闭式。
由闭式在 \(\Delta\to 0\) 的极限（或直接积分）得到 \(K(0)=1/6\)。
局部展开由 \(\sinh\) 与 \(\cosh\) 的 Taylor 展开直接给出。证毕。
\end{proof}

\begin{definition}[尺度相干 Gram 体积势]\label{def:xi-scale-coherence-gram-volume}
给定深度谱 \(\{\sigma_j\}_{j=1}^{\kappa}\) 与命题 \ref{prop:xi-scale-coherence-kernel-closed-form} 的核 \(K\)，定义尺度相干 Gram 矩阵
\[
G^{\mathrm{scale}}_{jk}:=K(\sigma_j-\sigma_k)\qquad(1\le j,k\le\kappa),
\]
并定义对应的尺度体积不变量与对数体积势
\[
V_{\mathrm{scale}}:=\sqrt{\det G^{\mathrm{scale}}},
\qquad
S_{\mathrm{scale}}:=-\log\det G^{\mathrm{scale}}.
\]
\end{definition}

\begin{proposition}[尺度碰撞的 Vandermonde 指数律：\texorpdfstring{$S_{\mathrm{scale}}$}{S_scale} 的普适发散阶恒为 \(2\)]\label{prop:xi-scale-coherence-collision-vandermonde}
设 \(\sigma_j=\sigma_\star+\varepsilon t_j\)（\(t_j\) 两两不同）且 \(\varepsilon\downarrow 0\)。
则存在常数 \(C_\kappa>0\)（仅依赖核 \(K\) 与 \(\kappa\)）使
\[
\det G^{\mathrm{scale}}(\boldsymbol{\sigma})
= C_\kappa\,\varepsilon^{\kappa(\kappa-1)}
\Bigl(\prod_{1\le j<k\le\kappa}(t_k-t_j)^2\Bigr)\cdot\bigl(1+o(1)\bigr).
\]
从而
\[
S_{\mathrm{scale}}(\boldsymbol{\sigma})
= -\log C_\kappa-\kappa(\kappa-1)\log\varepsilon-2\sum_{1\le j<k\le\kappa}\log|t_k-t_j|+o(1).
\]
因此每一对尺度碰撞 \(\sigma_j\to\sigma_k\) 都贡献严格的 \(2\log|\sigma_j-\sigma_k|^{-1}\) 发散阶。
\end{proposition}
\begin{proof}
由命题 \ref{prop:xi-scale-coherence-kernel-closed-form}，\(K\) 在 \(0\) 的邻域内为偶解析函数，故可写作
\[
K(\Delta)=\sum_{n\ge 0}\mu_{2n}\Delta^{2n},\qquad \mu_0=K(0)>0.
\]
将 \(\sigma_j=\sigma_\star+\varepsilon t_j\) 代入得
\(
G^{\mathrm{scale}}_{jk}=K(\varepsilon(t_j-t_k))
\)。
对矩阵元作 Taylor 展开并对行列式应用 confluent Vandermonde 型展开（等价地：把 \(\{K(\cdot-\sigma_j)\}\) 在 \(\varepsilon\to 0\) 下投影到由 \(\{\partial^r K(\cdot-\sigma_\star)\}_{r=0}^{\kappa-1}\) 张成的有限维子空间），得到
\[
\det G^{\mathrm{scale}}(\boldsymbol{\sigma})
=\varepsilon^{\kappa(\kappa-1)}\Bigl(\prod_{j<k}(t_k-t_j)^2\Bigr)\cdot\Bigl(\det H_\kappa+o(1)\Bigr),
\]
其中 \(H_\kappa\) 为由 \(\{\mu_{2n}\}\) 生成的 Hankel 型矩阵，\(\det H_\kappa>0\) 仅依赖于 \(K\) 与 \(\kappa\)。
令 \(C_\kappa:=\det H_\kappa\) 即得所示主项。
对数渐近由对两边取 \(-\log\) 并展开得到。证毕。
\end{proof}

\begin{proposition}[尺度相干核的 \texorpdfstring{$\zeta$}{zeta}-系数：\texorpdfstring{$K$}{K} 的 Taylor 系数被 \texorpdfstring{$\zeta(2n+2)$}{zeta(2n+2)} 完全封口]\label{prop:xi-scale-coherence-kernel-zeta-taylor}
记命题 \ref{prop:xi-scale-coherence-kernel-closed-form} 的偶展开 \(K(\Delta)=\sum_{n\ge 0}\mu_{2n}\Delta^{2n}\)。
则对一切 \(n\ge 0\) 有严格闭式
\[
\boxed{\;
\mu_{2n}
=\frac{K^{(2n)}(0)}{(2n)!}
=(-1)^n\,\frac{(2n+2)(2n+1)}{2^{2n+1}\pi^{2n+2}}\;\zeta(2n+2).\;}
\]
因此命题 \ref{prop:xi-scale-coherence-collision-vandermonde} 的常数 \(C_\kappa\) 可表示为由 \(\{\mu_{2n}\}\) 生成的 Hankel 型行列式（从而本质上为偶 \(\zeta\) 值的有理组合）。
\end{proposition}
\begin{proof}
由命题 \ref{prop:xi-bulk-curvature-tomography-identifiability} 的 Fourier 闭式，
\(
\widehat\phi(\xi)=\pi\xi/\sinh(\pi\xi)
\)。
又由卷积恒等式 \(K=\phi\ast\phi\)（\(\phi\) 为偶函数）得
\[
K(\Delta)=\frac{1}{2\pi}\int_{\RR}\widehat\phi(\xi)^2\,e^{\mathrm{i}\xi\Delta}\,\dd\xi
=\frac{1}{2\pi}\int_{\RR}\Bigl(\frac{\pi\xi}{\sinh(\pi\xi)}\Bigr)^{\!2}e^{\mathrm{i}\xi\Delta}\,\dd\xi.
\]
对 \(e^{\mathrm{i}\xi\Delta}\) 作幂级数展开并逐项积分（绝对收敛性由 \(\widehat\phi(\xi)=O(e^{-\pi|\xi|})\) 给出）得到
\[
\mu_{2n}
=\frac{1}{2\pi(2n)!}\int_{\RR}(\mathrm{i}\xi)^{2n}\Bigl(\frac{\pi\xi}{\sinh(\pi\xi)}\Bigr)^{\!2}\dd\xi
=\frac{(-1)^n\pi}{2(2n)!}\int_{\RR}\frac{\xi^{2n+2}}{\sinh^2(\pi\xi)}\,\dd\xi.
\]
利用级数展开
\(
\sinh^{-2}x=4\sum_{k\ge 1}k\,e^{-2kx}\ (x>0)
\)
并对偶函数化简为 \(2\int_0^\infty\)，得到
\[
\int_{\RR}\frac{\xi^{2n+2}}{\sinh^2(\pi\xi)}\,\dd\xi
=8\sum_{k\ge 1}k\int_0^\infty \xi^{2n+2}e^{-2\pi k\xi}\,\dd\xi
=8\sum_{k\ge 1}k\cdot\frac{(2n+2)!}{(2\pi k)^{2n+3}}
=\frac{(2n+2)!}{2^{2n}\pi^{2n+3}}\zeta(2n+2).
\]
代回并化简即得所示闭式。
关于 \(C_\kappa\) 的表述由命题 \ref{prop:xi-scale-coherence-collision-vandermonde} 证明中的 Hankel 结构与系数闭式立得。证毕。
\end{proof}

\begin{remark}[\texorpdfstring{$q=4$}{q=4} 的 \texorpdfstring{$\zeta$}{zeta} 共振起点：\texorpdfstring{$\zeta(4)$}{zeta(4)} 同时作为单点非高斯与两点相干的首个触发阶]\label{rem:xi-zeta4-resonance-threshold}
由命题 \ref{prop:xi-bulk-curvature-mgf-factorization-zeta-tower}，尺度曲率分布 \(p\) 的普适累积量在四阶首次引入偶 \(\zeta\) 常数：
\[
\kappa_{4}^{\mathrm{univ}}(p)=4!\,\frac{\zeta(4)}{2}=12\,\zeta(4)=\frac{2\pi^4}{15}.
\]
另一方面，由命题 \ref{prop:xi-scale-coherence-kernel-zeta-taylor} 的局部展开，
\[
\mu_{2}=-\,\frac{3}{2\pi^{4}}\zeta(4)=-\frac{1}{60},
\]
亦即两点尺度相干的首个非平凡几何修正（\(\Delta^2\) 项）同样由 \(\zeta(4)\) 统一封口。
因此在单点涨落与两点相干的双重层面上，\(\zeta(4)\) 构成该体系的首个共同算术触发阶。
\end{remark}

\begin{proposition}[深度谱的主化序与尺度曲率的凸序：更分散的边界必导致更强的尺度涨落]\label{prop:xi-depth-spectrum-majorization-convex-order}
设两组深度谱对数坐标
\(
\boldsymbol{\sigma}=(\sigma_1,\dots,\sigma_\kappa)
\)
与
\(
\boldsymbol{\sigma}'=(\sigma_1',\dots,\sigma_\kappa')
\)
满足同均值
\(
\kappa^{-1}\sum_j\sigma_j=\kappa^{-1}\sum_j\sigma_j'
\)，
且 \(\boldsymbol{\sigma}\) 主化（majorize）\(\boldsymbol{\sigma}'\)。
则对任意凸函数 \(F:\RR\to\RR\) 有
\[
\frac{1}{\kappa}\sum_{j=1}^{\kappa}F(\sigma_j)\ \ge\ \frac{1}{\kappa}\sum_{j=1}^{\kappa}F(\sigma_j').
\]
令 \(p_{\boldsymbol{\sigma}}(s):=\kappa^{-1}\sum_{j=1}^{\kappa}\phi(s-\sigma_j)\) 与 \(p_{\boldsymbol{\sigma}'}\) 为对应的尺度曲率分布。
则上述主化关系在凸序意义下被保序提升：对任意凸 \(F\) 都有
\[
\int_{\RR}F(s)\,p_{\boldsymbol{\sigma}}(s)\,\dd s\ \ge\ \int_{\RR}F(s)\,p_{\boldsymbol{\sigma}'}(s)\,\dd s.
\]
特别地，所有中心偶矩满足单调性：深度谱越分散，则尺度曲率的方差、峰度及更高偶阶中心矩必不减。
\end{proposition}
\begin{proof}
主化蕴含的 Jensen 型不等式为 Hardy--Littlewood--P\'olya 理论的标准结论：对任意凸 \(F\) 有离散平均单调性。
另一方面，令随机变量 \(U\) 在 \(\{1,\dots,\kappa\}\) 上均匀，取 \(S:=\sigma_U\)、\(S':=\sigma'_U\)。
再令 \(Z\) 独立且密度为 \(\phi\)，则 \(S+Z\) 与 \(S'+Z\) 的密度分别为 \(p_{\boldsymbol{\sigma}}\) 与 \(p_{\boldsymbol{\sigma}'}\)。
由凸序的平移不变性（独立和保持凸序）得对任意凸 \(F\)，
\(
\EE[F(S+Z)]\ge \EE[F(S'+Z)]
\)，
这正是所示积分不等式。
取 \(F(s)=(s-\EE[s])^{2n}\) 得中心偶矩单调性。证毕。
\end{proof}

% (User input window)
%
% 在 logistic 核层析口径下，曲率概率密度 \(p\) 的反演不仅是可识别的，而且具有全显式算子表达；
% 同一乘子结构同时强制噪声放大、零点密度与比特复杂度的普适阈值。

命题 \ref{prop:xi-bulk-curvature-logistic-kernel-decomposition}--\ref{prop:xi-bulk-curvature-tomography-identifiability} 给出曲率核的有限原子分解与可识别性；
在同一记号下，可进一步将反演写成显式算子，并由同一乘子结构推出稳定性、零点密度与精度预算的普适封口。

\begin{proposition}[尺度曲率的精确反卷积：logistic 平滑的显式逆算子]\label{prop:xi-bulk-curvature-exact-deconvolution-operator}
\leanverified{paper\_xi\_bulk\_curvature\_exact\_deconvolution\_operator}
在推论 \ref{cor:xi-bulk-curvature-probability-moments} 的设定下，记
\[
\nu:=\frac{1}{\kappa}\sum_{j=1}^{\kappa}\delta_{\sigma_j},
\qquad
\phi(s)=\frac{e^{s}}{(1+e^{s})^{2}}=\frac{1}{4\cosh^{2}(s/2)},
\qquad
p=\phi*\nu.
\]
令 \(\mathcal D\) 为 \(\mathcal S'(\RR)\) 上的 Fourier 乘子算子（在 \(\xi=0\) 处取连续延拓值）
\[
\widehat{(\mathcal D f)}(\xi):=\frac{\sinh(\pi\xi)}{\pi\xi}\,\widehat f(\xi).
\]
则在温和分布意义下有严格恒等式
\[
\boxed{\;\mathcal D p=\nu.\;}
\]
并且 \(\mathcal D\) 具有 Weierstrass 乘积表示（作为 Fourier 乘子算子作用于 \(\mathcal S(\RR)\)）
\[
\boxed{\;
\mathcal D=\prod_{n\ge1}\Bigl(1-\frac{1}{n^{2}}\partial_{s}^{2}\Bigr)
=\frac{\sinh\!\bigl(\pi\sqrt{-\partial_{s}^{2}}\bigr)}{\pi\sqrt{-\partial_{s}^{2}}}\; }.
\]
因此“尺度平滑 \(p=\phi*\nu\)”在该体系内并非不可逆噪声通道，而是由一条全显式的整函数型无限阶算子精确可逆。
\end{proposition}
\begin{proof}
由命题 \ref{prop:xi-bulk-curvature-tomography-identifiability} 中的闭式 Fourier 变换
\(
\widehat\phi(\xi)=\pi\xi/\sinh(\pi\xi)
\)
与卷积定理，得
\(
\widehat p(\xi)=\widehat\phi(\xi)\widehat\nu(\xi)
\)。
两边同乘 \(\sinh(\pi\xi)/(\pi\xi)\) 即得
\(
\widehat\nu(\xi)=\frac{\sinh(\pi\xi)}{\pi\xi}\widehat p(\xi)
=\widehat{(\mathcal D p)}(\xi)
\)，从而 \(\mathcal Dp=\nu\)。
乘积表示由经典整函数恒等式
\(
\sinh(\pi z)/(\pi z)=\prod_{n\ge1}(1+z^{2}/n^{2})
\)
与 Fourier 侧替换 \(z^{2}\mapsto \xi^{2}\)（等价于物理侧 \(1+z^{2}/n^{2}\mapsto 1-\partial_{s}^{2}/n^{2}\)）得到。证毕。
\end{proof}

% (User input window)

\begin{corollary}[尺度形状因子的模平方全息：深度差分多重集的加权计数]\label{cor:xi-bulk-curvature-modsquare-holography}
\leanverified{paper\_xi\_bulk\_curvature\_modsquare\_holography}
令
\(
\widehat p(\xi)=\int_{\RR}e^{-i\xi s}p(s)\,\dd s
\)。
则有严格恒等式
\[
\boxed{\;
\Bigl|\frac{\sinh(\pi\xi)}{\pi\xi}\widehat p(\xi)\Bigr|^2
=|\widehat\nu(\xi)|^2
=\frac{1}{\kappa^2}\sum_{j,k=1}^{\kappa}e^{-i\xi(\sigma_j-\sigma_k)}
=\frac{1}{\kappa^2}\sum_{\Delta\in\mathcal D}m(\Delta)e^{-i\xi\Delta}\; }.
\]
其中 \(\mathcal D:=\{\sigma_j-\sigma_k\}_{j,k}\) 为差分多重集，\(m(\Delta)\) 为其重数。
因此在剥离通用乘子 \(\widehat\phi(\xi)=\pi\xi/\sinh(\pi\xi)\) 后，模平方数据 \(|\widehat p(\xi)|^2\) 等价于深度差分谱的特征函数，从而二体几何可在不恢复相位的条件下被直接审计。
\end{corollary}
\begin{proof}
由命题 \ref{prop:xi-bulk-curvature-exact-deconvolution-operator} 得
\(
\widehat\nu(\xi)=\frac{\sinh(\pi\xi)}{\pi\xi}\widehat p(\xi)
\)。
取模平方并用
\(
\widehat\nu(\xi)=\kappa^{-1}\sum_{j=1}^{\kappa}e^{-i\xi\sigma_j}
\)
展开即得。证毕。
\end{proof}

\begin{corollary}[一般位置下的深度谱唯一性：由差分多重集刚性反演点集]\label{cor:xi-bulk-curvature-difference-multiset-rigid-inversion}
\leanverified{paper\_xi\_bulk\_curvature\_difference\_multiset\_rigid\_inversion}
设深度谱 \(\{\sigma_j\}_{j=1}^{\kappa}\subset\RR\) 满足一般位置条件：除 \(0\) 外，差分多重集 \(\mathcal D=\{\sigma_j-\sigma_k\}_{j,k}\) 不含重元，即对任意 \(\Delta\neq 0\) 有 \(m(\Delta)\in\{0,1\}\)。
则差分多重集 \((\mathcal D,m)\) 唯一确定无序多重集 \(\{\sigma_j\}\) 于整体平移与反射不变换之下；并且该反演可由有限步组合算法实现。
结合推论 \ref{cor:xi-bulk-curvature-modsquare-holography}，在该一般位置假设下，模平方数据 \(|\widehat p(\xi)|^{2}\) 已足以将深度谱锁定到仅剩 \(\RR\rtimes\ZZ_2\) 的全局对称自由度。
\end{corollary}
\begin{proof}[证明要点]
这是经典 turnpike/Golomb ruler 反演问题的标准刚性结论：最大差分给出两端点间距，递归地从差分多重集中剥离与端点兼容的差分，并在一般位置假设下消除分支，从而得到唯一解；平移与反射为唯一残余对称。证毕。
\end{proof}

\begin{theorem}[Hankel--Prony 学习的零点谱稳定性与可递归误差证书]\label{thm:xi-hankel-prony-root-stability-recursive-certificate}
\leanverified{paper\_xi\_hankel\_prony\_root\_stability\_recursive\_certificate}
设 \(\zeta_1,\dots,\zeta_r\in\CC^\times\) 两两不同，\(c_1,\dots,c_r\in\CC\setminus\{0\}\)，并定义矩序列
\[
m_k:=\sum_{j=1}^{r}c_j\zeta_j^k,\qquad k=0,1,\dots,2r-1.
\]
令 Hankel 矩阵
\[
H:=(m_{i+j})_{0\le i,j\le r-1}\in\Mat_r(\CC),
\]
并假设 \(H\) 可逆，记 \(\sigma_{\min}:=\sigma_{\min}(H)>0\)。
给定噪声观测 \(\tilde m_k=m_k+\varepsilon_k\) 且 \(|\varepsilon_k|\le \epsilon\)，构造 \(\tilde H=(\tilde m_{i+j})_{0\le i,j\le r-1}\)。
则存在常数
\[
C=C\!\left(r,\ \max_j|c_j|,\ \max_j|\zeta_j|\right)>0
\]
使得当 \(\epsilon\le \sigma_{\min}/(2C)\) 时，由 Prony--Hankel 过程从 \(\{\tilde m_k\}_{k=0}^{2r-1}\) 所恢复的估计根 \(\tilde\zeta_1,\dots,\tilde\zeta_r\) 可重标记满足
\[
\max_{1\le j\le r}|\tilde\zeta_j-\zeta_j|\ \le\ \frac{C}{\sigma_{\min}}\,\epsilon.
\]
并且当追加矩观测 \(\tilde m_{2r},\tilde m_{2r+1},\dots\) 以形成更高阶 Hankel 扩维矩阵时，\(\sigma_{\min}\) 的下界与上述上界可由 Schur 补与最小奇异值扰动估计逐步更新，从而在每一阶生成后验可核验的稳定性证书。
\end{theorem}
\begin{proof}[证明要点]
矩分解可写为 \(H=VCV^{\mathsf T}\)，其中 \(V=(\zeta_j^i)_{0\le i\le r-1,\,1\le j\le r}\) 为 Vandermonde 矩阵且 \(C=\mathrm{diag}(c_1,\dots,c_r)\)。
由 \(|\varepsilon_k|\le\epsilon\) 得 \(\|\tilde H-H\|_{\mathrm{op}}\le C_0(r)\epsilon\)。
当 \(\epsilon\le \sigma_{\min}/(2C)\) 时，可保证 \(\tilde H\) 可逆且 \(\|\tilde H^{-1}\|_{\mathrm{op}}\le 2/\sigma_{\min}\)。
Prony 重构可表述为从 \(\tilde H\) 解出线性递推（或等价 companion 矩阵）并取其特征根；该映射在根互异域内对数据是局部 Lipschitz。
将线性系统的条件数界（由 \(\sigma_{\min}^{-1}\) 控制）与多项式根对系数的稳定性界（例如由 companion 矩阵的 Bauer--Fike 型估计）合并，即得所示 \(\frac{C}{\sigma_{\min}}\epsilon\) 上界；常数 \(C\) 仅依赖于 \(r\) 与 \(\max_j|c_j|\)、\(\max_j|\zeta_j|\) 的上界。
递归更新部分由 Hankel 扩维的块矩阵 Schur 补公式与 Weyl 型奇异值扰动界给出。证毕。
\end{proof}

\begin{proposition}[累积矩噪声下的 Hankel 反演稳定性：逆矩阵差分与特征值扰动界]\label{prop:xi-hankel-inversion-noise-stability-neumann-bauerfike}
\leanverified{paper\_xi\_hankel\_inversion\_noise\_stability\_neumann\_bauerfike}
令 \(H_0,H_1\in\Mat_r(\CC)\)，并设 \(H_0\) 可逆。
给定扰动 \(\widetilde H_0,\widetilde H_1\in\Mat_r(\CC)\)，假设
\begin{equation}\label{eq:xi-hankel-neumann-smallness}
\boxed{\;
\|H_0^{-1}\|\,\|\widetilde H_0-H_0\|\ \le\ \frac12\;}
\end{equation}
其中 \(\|\cdot\|\) 为任意次乘法矩阵范数（例如算子范数）。
则：
\begin{enumerate}
  \item \(\widetilde H_0\) 可逆且满足
  \[
  \boxed{\;
  \|\widetilde H_0^{-1}-H_0^{-1}\|
  \ \le\
  2\,\|H_0^{-1}\|^{2}\,\|\widetilde H_0-H_0\|\;}
  \]
  并且 \(\|\widetilde H_0^{-1}\|\le 2\|H_0^{-1}\|\)；
  \item 令
  \[
  S:=H_0^{-1}H_1,\qquad \widetilde S:=\widetilde H_0^{-1}\widetilde H_1,
  \]
  则有显式扰动界
  \[
  \boxed{\;
  \|\widetilde S-S\|
  \ \le\
  4\,\|H_0^{-1}\|\Bigl(\|\widetilde H_1-H_1\|+\|S\|\,\|\widetilde H_0-H_0\|\Bigr)\;}
  \]
  \item 若 \(S\) 可对角化，即 \(S=PDP^{-1}\)（\(D\) 对角），则对任意 \(\widetilde S\) 的特征值 \(\widetilde\xi\) 有
  \[
  \boxed{\;
  \mathrm{dist}\bigl(\widetilde\xi,\sigma(S)\bigr)\ \le\ \kappa(P)\,\|\widetilde S-S\|,
  \qquad \kappa(P):=\|P\|\,\|P^{-1}\|\;}
  \]
  因而在零点重构 \(\widetilde\zeta:=1/\widetilde\xi\) 的口径下，若 \(\sigma(S)\) 与 \(0\) 保持正距离，则零点集合在 Hausdorff 距离意义下具有显式 Lipschitz 稳定性常数。
\end{enumerate}
\end{proposition}
\begin{proof}
(1) 写作 \(\widetilde H_0=H_0(I+E)\)，其中 \(E:=H_0^{-1}(\widetilde H_0-H_0)\)。
由 \eqref{eq:xi-hankel-neumann-smallness} 得 \(\|E\|\le 1/2\)，故 \(I+E\) 可逆且
\(\|(I+E)^{-1}\|\le \sum_{k\ge 0}\|E\|^k\le 2\)。
于是 \(\widetilde H_0^{-1}=(I+E)^{-1}H_0^{-1}\) 且 \(\|\widetilde H_0^{-1}\|\le 2\|H_0^{-1}\|\)。
并且由恒等式
\(
\widetilde H_0^{-1}-H_0^{-1}=\widetilde H_0^{-1}(H_0-\widetilde H_0)H_0^{-1}
\)
得
\[
\|\widetilde H_0^{-1}-H_0^{-1}\|
\le \|\widetilde H_0^{-1}\|\,\|\widetilde H_0-H_0\|\,\|H_0^{-1}\|
\le 2\|H_0^{-1}\|^2\|\widetilde H_0-H_0\|.
\]
(2) 展开
\[
\widetilde S-S
=(\widetilde H_0^{-1}-H_0^{-1})H_1+\widetilde H_0^{-1}(\widetilde H_1-H_1),
\]
并用 \(\|H_1\|\le \|H_0\|\,\|S\|\) 与 \(\|\widetilde H_0^{-1}\|\le 2\|H_0^{-1}\|\)，结合 (1) 的界即可得到所示估计（常数用粗估吸收为 \(4\)）。
(3) 由 Bauer--Fike 定理立得特征值扰动界。若 \(\inf_{\xi\in\sigma(S)}|\xi|>0\)，则映射 \(\xi\mapsto 1/\xi\) 在该谱邻域上 Lipschitz，从而零点重构的 Hausdorff 稳定性随之成立。证毕。
\end{proof}

\begin{corollary}[超分辨率屏障常数 \texorpdfstring{$\pi$}{pi}：深度分离的稳定阈值为 \texorpdfstring{$\exp(-\pi/\Delta)$}{exp(-pi/Delta)}]\label{cor:xi-bulk-curvature-superresolution-pi-barrier}
设观测到的曲率概率密度为 \(\widetilde p=p+\eta\)，并假定 Fourier 噪声在带宽 \(|\xi|\le\Omega\) 内满足
\[
\sup_{|\xi|\le\Omega}|\widehat\eta(\xi)|\le\varepsilon.
\]
令线性反卷积输出为 \(\widetilde\nu:=\mathcal D\widetilde p\)。
由命题 \ref{prop:xi-bulk-curvature-exact-deconvolution-operator} 的反卷积乘子可知，在最坏情形噪声下，
\[
\sup_{|\xi|\le\Omega}\bigl|\widehat{(\widetilde\nu-\nu)}(\xi)\bigr|
\ \gtrsim\
\varepsilon\cdot \frac{\sinh(\pi\Omega)}{\pi\Omega}
\ \asymp\
\varepsilon\,e^{\pi\Omega}.
\]
因此若希望在尺度轴上以分离度 \(\Delta\) 稳定地区分原子（通常需要使用 \(\Omega\gtrsim 1/\Delta\) 的频率），则噪声预算必须在量级上满足
\[
\varepsilon \lesssim \exp\!\Bigl(-\frac{\pi}{\Delta}\Bigr),
\qquad\text{等价地}\qquad
\Delta \gtrsim \frac{\pi}{\log(1/\varepsilon)}.
\]
该阈值不依赖于 \(\kappa\) 与 \(\{\sigma_j\}\) 的几何细节，常数 \(\pi\) 为普适屏障常数。
\end{corollary}
\begin{proof}[证明要点]
由 \(\widetilde\nu-\nu=\mathcal D\eta\) 与乘子定义，
\(
|\widehat{(\widetilde\nu-\nu)}(\xi)|
=\frac{\sinh(\pi|\xi|)}{\pi|\xi|}\,|\widehat\eta(\xi)|
\)。
在噪声约束下选择相位对齐的 \(\widehat\eta\) 可达到最坏情形放大因子 \(\sinh(\pi\Omega)/(\pi\Omega)\)。
再与分辨率--带宽互反标尺 \(\Omega\gtrsim 1/\Delta\) 合并即可。证毕。
\end{proof}

% (User input window)

\begin{remark}[双重不可逾越界：原子数—分离度—精度的联合极限]\label{rem:xi-bulk-curvature-double-barrier-kappa-separation-precision}
在深度谱恢复问题中，分离度 \(\Delta_{\min}\) 同时受到两条互不替代的制度性门槛约束：
\begin{itemize}
\item \textbf{体积稳定门槛（Gram 侧）}：
若希望 Hellinger--Gram 体积不发生簇化塌缩，则需进入准正交相的分离阈值（推论 \ref{cor:xi-hellinger-quasi-orthogonal-phase}）。
反之，一旦 \(\Delta_{\min}\) 被迫随 \(\kappa\) 缩小，则簇化条件数以 \(\kappa(\kappa-1)\) 次幂制度爆炸（推论 \ref{cor:xi-hellinger-gram-cluster-conditioning-exponent}）。
\item \textbf{可辨识门槛（反卷积侧）}：
要在分离度 \(\Delta_{\min}\) 下稳定区分原子，噪声预算必须满足
\(
\varepsilon\lesssim \exp(-\pi/\Delta_{\min})
\)
（推论 \ref{cor:xi-bulk-curvature-superresolution-pi-barrier}）。
在固定跨度 \(L=\sigma_\kappa-\sigma_1\) 的约束下有 \(\Delta_{\min}\le L/(\kappa-1)\)，从而噪声门槛进一步压缩为
\[
\varepsilon\lesssim \exp\!\Bigl(-\frac{\pi(\kappa-1)}{L}\Bigr),
\]
即精度成本对 \(\kappa\) 呈指数级增长。
\end{itemize}
因此在任何限制 \(\Delta_{\min}\) 随 \(\kappa\) 缩小的预算口径下，原子数提升与稳定恢复不可同时进入任意高秩与任意小误差的双极限；可行域由 \(\Delta_{\min}\) 的双重门槛共同截断。
\end{remark}

\begin{proposition}[极点--整函数二分解：尺度曲率母函数的通用晶格奇性与深度因子]\label{prop:xi-bulk-curvature-moment-generating-factorization}
\leanverified{paper\_xi\_bulk\_curvature\_moment\_generating\_factorization}
定义母函数（双侧 Laplace 变换）
\[
M_p(t):=\int_{\RR}e^{ts}p(s)\,\dd s,
\qquad |\Re t|<1.
\]
则存在严格因子化
\[
\boxed{\;
M_p(t)=\frac{\pi t}{\sin(\pi t)}\,E(t),\qquad
E(t):=\int_{\RR}e^{ts}\,\dd\nu(s)=\frac{1}{\kappa}\sum_{j=1}^{\kappa}e^{t\sigma_j}.\;}
\]
其中 \(\frac{\pi t}{\sin(\pi t)}\) 给出全部通用奇性：在每个非零整数 \(n\in\ZZ\setminus\{0\}\) 处具有简单极点；
而 \(E\) 为整函数并携带全部深度信息。
特别地，
\[
\boxed{\;\frac{\sin(\pi t)}{\pi t}\,M_p(t)=E(t)\;}
\]
是有限指数和，构成深度谱的解析指纹。
\end{proposition}
\begin{proof}
由 \(p=\phi*\nu\) 与 Fubini 定理（\(|\Re t|<1\) 保障绝对可积），得
\[
M_p(t)=\int_{\RR}\int_{\RR}e^{t(s+u)}\phi(s)\,\dd s\,\dd\nu(u)
=\Bigl(\int_{\RR}e^{ts}\phi(s)\,\dd s\Bigr)\cdot\Bigl(\int_{\RR}e^{tu}\,\dd\nu(u)\Bigr).
\]
对 \(\phi(s)=e^{s}(1+e^{s})^{-2}\)，作代换 \(x=e^{s}\)（\(\dd s=\dd x/x\)）得
\[
\int_{\RR}e^{ts}\phi(s)\,\dd s
=\int_{0}^{\infty}\frac{x^{t}}{(1+x)^{2}}\,\dd x
=\mathrm{B}(t+1,1-t)
=\Gamma(t+1)\Gamma(1-t)
=\frac{\pi t}{\sin(\pi t)}.
\]
从而得到所示因子化。证毕。
\end{proof}

\begin{remark}[尺度 \texorpdfstring{$\zeta$}{zeta} 接口的记号：母函数即尺度 \texorpdfstring{$\zeta$}{zeta} 函数]\label{rem:xi-bulk-curvature-scale-zeta-notation}
母函数 \(M_p(t)\) 亦可记为尺度 \(\zeta\) 接口
\[
\zeta_{\mathrm{scale}}(t):=\int_{\RR}e^{ts}p(s)\,\dd s=M_p(t),
\qquad
|\Re t|<1.
\]
由命题 \ref{prop:xi-bulk-curvature-moment-generating-factorization}，\(\zeta_{\mathrm{scale}}\) 具有唯一的亚纯延拓；
其全部整数晶格极点由通用因子 \(\pi t/\sin(\pi t)\) 承载，而深度信息完全由整函数 \(E(t)\) 承载。
\end{remark}

\begin{corollary}[整数极点留数的深度 Dirichlet 数据：有限个局部数据的代数重建]\label{cor:xi-bulk-curvature-residue-dirichlet-data}
在命题 \ref{prop:xi-bulk-curvature-moment-generating-factorization} 的设定下，对任意 \(n\in\NN\) 有显式留数公式
\[
\boxed{\;
\operatorname*{Res}_{t=n}M_p(t)=(-1)^n\,n\,E(n),\qquad
\operatorname*{Res}_{t=-n}M_p(t)=(-1)^n\,(-n)\,E(-n).\;}
\]
令 \(y_j:=e^{\sigma_j}=\delta_j^{-1}\)。
则
\[
E(n)=\frac{1}{\kappa}\sum_{j=1}^{\kappa}y_j^{n}.
\]
若 \(\{y_j\}\) 两两不同，则有限序列 \(\{E(n)\}_{n=1}^{2\kappa}\)（等价地，\(\{\operatorname*{Res}_{t=\pm n}M_p(t)\}_{n=1}^{2\kappa}\)）唯一确定无序多重集 \(\{y_j\}\)，从而唯一确定 \(\{\delta_j\}\)。
\end{corollary}
\begin{proof}
由
\(
\frac{\pi t}{\sin(\pi t)}
\)
在 \(t=n\) 的留数为 \((-1)^n n\)，与因子化式相乘即得留数公式。
第二部分为经典幂和反演：\(\{y_j\}\) 两两不同下，幂和序列满足阶数为 \(\kappa\) 的线性递推；
由 \(\{E(n)\}_{n=1}^{2\kappa}\) 可构造 Hankel 系统解出递推系数，从而得到以 \(\{y_j\}\) 为根的首一多项式并恢复 \(\{y_j\}\)，进而恢复 \(\{\delta_j=y_j^{-1}\}\)。证毕。
\end{proof}

% (User input window)

\begin{corollary}[整数留数的极值重建：最大深度与重数的无损读出]\label{cor:xi-bulk-curvature-residue-extremal-recovery}
在推论 \ref{cor:xi-bulk-curvature-residue-dirichlet-data} 的设定下，令
\[
R_n:=(-1)^n\,\frac{1}{n}\operatorname*{Res}_{t=n}\zeta_{\mathrm{scale}}(t)=E(n)\qquad(n\in\NN).
\]
记最大深度与其重数为
\[
\sigma_{\max}:=\max_{1\le j\le\kappa}\sigma_j,
\qquad
r_{\max}:=\card\{j:\ \sigma_j=\sigma_{\max}\}.
\]
则有严格极限公式
\[
\boxed{\;
\sigma_{\max}=\lim_{n\to\infty}\frac{1}{n}\log\bigl(\kappa R_n\bigr),\qquad
r_{\max}=\lim_{n\to\infty}\kappa R_n\,e^{-n\sigma_{\max}}\in\NN.\;}
\]
因此“最深缺陷”及其重数无需全谱反演即可由整数极点的留数序列直接证书化。
\end{corollary}
\begin{proof}[证明要点]
由 \(R_n=E(n)=\kappa^{-1}\sum_{j=1}^{\kappa}e^{n\sigma_j}\)，将最大深度项分离得
\[
\kappa R_n=r_{\max}e^{n\sigma_{\max}}+\sum_{\sigma_j<\sigma_{\max}}e^{n\sigma_j}
=r_{\max}e^{n\sigma_{\max}}\Bigl(1+O(e^{-n(\sigma_{\max}-\sigma_2)})\Bigr),
\]
其中 \(\sigma_2\) 为次大深度（若不存在则误差视为 \(0\)）。
两端取对数并除以 \(n\) 即得 \(\sigma_{\max}\) 的极限；再将上式乘以 \(e^{-n\sigma_{\max}}\) 取极限即得 \(r_{\max}\)。证毕。
\end{proof}

\begin{remark}[尺度 Euler--Möbius 原语账本：幂和矩与 primitive 序列的互换]\label{rem:xi-bulk-curvature-euler-mobius-ledger}
令 \(y_j:=e^{\sigma_j}=\delta_j^{-1}\)，并记幂和
\(
p_n:=\sum_{j=1}^{\kappa}y_j^{n}=\kappa E(n)
\)（\(n\in\NN\)）。
引入形式幂级数
\[
\mathfrak Z_{\mathrm{scale}}(z)
:=\prod_{j=1}^{\kappa}\frac{1}{1-y_j z}
=\exp\Bigl(\sum_{n\ge 1}\frac{p_n}{n}z^{n}\Bigr)
\qquad(\text{形式意义}).
\]
并定义其 primitive 系数（Möbius 反演）
\[
\pi^{\mathrm{prim}}_n
:=\frac{1}{n}\sum_{d\mid n}\mu(d)\,p_{n/d}
=\frac{1}{n}\sum_{d\mid n}\mu(d)\sum_{j=1}^{\kappa}y_j^{n/d}.
\]
则有等价 Euler 分解
\[
\mathfrak Z_{\mathrm{scale}}(z)=\prod_{n\ge 1}\frac{1}{(1-z^{n})^{\pi^{\mathrm{prim}}_n}}
\qquad(\text{形式意义}),
\]
从而 \(\{E(n)\}\)（或 \(\{p_n\}\)）给出“总量矩”账本，而 \(\{\pi^{\mathrm{prim}}_n\}\) 给出其 Euler--Möbius 的原语分解；二者互相可由有限步代数运算转写。
\end{remark}

\begin{proposition}[原语账本的指数主导律：最大深度锁定 primitive 指数率]\label{prop:xi-bulk-curvature-primitive-max-depth-dominance}
沿用注 \ref{rem:xi-bulk-curvature-euler-mobius-ledger} 的记号，并设
\[
\Lambda:=\max_{1\le j\le\kappa}y_j,\qquad
r:=\card{\{j:\ y_j=\Lambda\}},
\qquad
\Lambda_2:=\max\{y_j:\ y_j<\Lambda\}\ \ (\text{若存在}).
\]
若 \(\Lambda>1\)，则对 \(n\to\infty\) 有
\[
\boxed{\;
\pi^{\mathrm{prim}}_n=\frac{r}{n}\Lambda^n
+O\!\Bigl(\frac{1}{n}\Lambda^{n/2}\Bigr)
+O\!\Bigl(\frac{1}{n}\Lambda_2^{\,n}\Bigr)\;}
\]
（第二项仅在 \(n\) 含非平凡因子时出现，且以 \(\Lambda^{n/2}\) 为最坏情形量级）。
从而指数率由最大深度唯一锁定：
\[
\boxed{\;\lim_{n\to\infty}\frac{1}{n}\log\bigl|n\pi^{\mathrm{prim}}_n\bigr|=\log\Lambda.\;}
\]
\end{proposition}
\begin{proof}
由定义 \(n\pi^{\mathrm{prim}}_n=\sum_{d\mid n}\mu(d)\,p_{n/d}\) 且 \(p_m=\sum_{j}y_j^m\)，得
\[
p_n=r\Lambda^n+O(\Lambda_2^{\,n}),
\]
其中 \(\Lambda_2\) 不存在时可将 \(O(\Lambda_2^{\,n})\) 理解为 \(0\)。
另一方面，对任意 \(d\mid n\) 且 \(d\ge 2\) 有 \(n/d\le n/2\)，从而
\[
|p_{n/d}|\le \sum_{j=1}^{\kappa}y_j^{n/d}\le \kappa\,\Lambda^{n/d}\le \kappa\,\Lambda^{n/2}.
\]
将 \(d=1\) 项与其余项分离即可得到所示展开。
最后，由 \(\Lambda>1\) 可知 \(\Lambda^{n/2}=o(\Lambda^n)\) 且 \(\Lambda_2^{\,n}=o(\Lambda^n)\)，故 \(n\pi^{\mathrm{prim}}_n\) 的指数增长率为 \(\log\Lambda\)。证毕。
\end{proof}

\begin{remark}[隐藏圆谱几何：反卷积乘子与圆 Laplacian 的 \texorpdfstring{$\zeta$}{zeta}-行列式比]\label{rem:xi-bulk-curvature-circle-zeta-determinant}
令 \(L=-\partial_\theta^2\) 为圆 \(S^1\) 上 Laplacian，其正谱为 \(\{n^2\}_{n\ge1}\)。
其 \(\zeta\)-正规化行列式比满足
\[
\frac{\det_\zeta(L+\xi^{2})}{\det_\zeta(L)}
=\prod_{n\ge1}\Bigl(1+\frac{\xi^{2}}{n^{2}}\Bigr)
=\frac{\sinh(\pi|\xi|)}{\pi|\xi|}.
\]
因此命题 \ref{prop:xi-bulk-curvature-exact-deconvolution-operator} 的反卷积乘子精确等同于圆 Laplacian 的 \(\zeta\)-行列式比；
常数 \(\pi\) 的出现对应于该隐藏圆谱几何在尺度轴上的投影。
\end{remark}

\begin{proposition}[尺度相干核与尺度体积势：短程 Vandermonde 排斥与长程指数正交化]\label{prop:xi-bulk-curvature-scale-gram-two-phase}
在推论 \ref{cor:xi-bulk-curvature-probability-moments} 的设定下，定义尺度相干核
\[
K(\Delta):=\int_{\RR}\phi(s)\phi(s-\Delta)\,\dd s=(\phi*\phi)(\Delta).
\]
则 \(K\) 为偶函数且具有闭式
\[
\boxed{\;
K(\Delta)=\frac{|\Delta|\coth\!\bigl(|\Delta|/2\bigr)-2}{4\sinh^{2}\!\bigl(|\Delta|/2\bigr)}\; } ,
\qquad
K(0)=\frac{1}{6}.
\]
对尺度谱 \(\{\sigma_j\}_{j=1}^{\kappa}\) 定义 Gram 矩阵
\(
G^{\mathrm{scale}}_{jk}:=K(\sigma_j-\sigma_k)
\)
与尺度体积势
\(
S_{\mathrm{scale}}:=-\log\det G^{\mathrm{scale}}
\)。
则存在两端极限结构：
当最小分离 \(\Delta_{\min}:=\min_{j\neq k}|\sigma_j-\sigma_k|\to\infty\) 时，
\[
\det G^{\mathrm{scale}}=(K(0))^{\kappa}\bigl(1+o(1)\bigr)=\Bigl(\frac16\Bigr)^{\kappa}\bigl(1+o(1)\bigr),
\qquad
S_{\mathrm{scale}}=\kappa\log 6+o(1),
\]
且核尾部满足 \(K(\Delta)\sim (|\Delta|-2)e^{-|\Delta|}\)，从而非对角耦合以 \(e^{-|\Delta|}\) 指数湮灭。
另一方面，当 \(\sigma_j=\sigma_\star+\varepsilon t_j\) 且 \(\varepsilon\downarrow0\)（\(t_j\) 两两不同）时，
\[
S_{\mathrm{scale}}
=\kappa(\kappa-1)\log\frac{1}{\varepsilon}
+2\sum_{1\le j<k\le \kappa}\log\frac{1}{|t_k-t_j|}
+O(1),
\]
表现出 Vandermonde 型对数排斥。
\end{proposition}
\begin{proof}
闭式由直接积分计算得到：对 \(\phi(s)=\frac{1}{4\cosh^{2}(s/2)}\) 令 \(s=2u\)，化为
\(
K(\Delta)=\frac18\int_{\RR}\frac{1}{\cosh^{2}(u)\cosh^{2}(u-\Delta/2)}\,\dd u
\)，再用双曲函数恒等式积分可得所示表达；\(K(0)=\int \phi^2=\frac16\) 由 \(\int_{\RR}\cosh(u)^{-4}\dd u=\frac43\) 得到。
\par
长程相由 \(G^{\mathrm{scale}}=K(0)I+O(\max_{j\neq k}K(\sigma_j-\sigma_k))\) 与行列式连续性推出；尾项渐近由闭式在 \(|\Delta|\to\infty\) 时展开得到。
\par
短程相可由平移族 \(\phi(\cdot-\sigma_\star-\varepsilon t_j)\) 在 \(L^2\) 中对 \(\varepsilon\) 的 Taylor 展开得到：该族在 \(\varepsilon\to0\) 时收敛到导数塔 \(\{\phi^{(r)}(\cdot-\sigma_\star)\}_{r=0}^{\kappa-1}\) 的张成；
Gram 行列式的首阶非零项由 confluent Vandermonde 结构给出，因而出现 \(\varepsilon^{\kappa(\kappa-1)}\prod_{j<k}(t_k-t_j)^2\) 的尺度律。
取对数即得所示展开。证毕。
\end{proof}

\begin{proposition}[深度跨度的解析指纹：去极点后的整函数零点密度给出跨度]\label{prop:xi-bulk-curvature-depth-span-zero-density}
在命题 \ref{prop:xi-bulk-curvature-moment-generating-factorization} 的设定下，记
\[
E(t)=\frac{\sin(\pi t)}{\pi t}\,M_p(t)=\frac{1}{\kappa}\sum_{j=1}^{\kappa}e^{t\sigma_j}.
\]
令深度跨度
\(
L:=\max_j\sigma_j-\min_j\sigma_j
\)。
设 \(N_E(R)\) 为 \(E\) 在圆盘 \(\{t\in\CC:\ |t|\le R\}\) 内的零点计数函数（按重数计），则存在极限
\[
\boxed{\;\lim_{R\to\infty}\frac{N_E(R)}{R}=\frac{L}{\pi}.\;}
\]
因此 \(L\) 可由母函数 \(M_p\) 去除通用晶格因子后、仅凭整函数 \(E\) 的零点密度无损读出。
\end{proposition}
\begin{proof}[证明要点]
因子 \(e^{t\min_j\sigma_j}\) 不改变零点集合，可将 \(E\) 化为形如 \(\sum_{j}c_j e^{a_j t}\) 的指数多项式（\(0=a_1<\cdots<a_\kappa=L\)）。
该类函数属 Cartwright 型整函数，其指示图为实轴上的线段 \([0,L]\)，由 Cartwright 密度定理可得零点计数满足 \(N_E(R)\sim (L/\pi)R\)。证毕。
\end{proof}

\begin{corollary}[稳定反演的高斯正则化：最小可辨尺度的普适门槛]\label{cor:xi-bulk-curvature-gaussian-regularization-threshold}
引入高斯正则化 \(\nu_\sigma:=G_\sigma*\nu\)（Fourier 乘子 \(e^{-\sigma^{2}\xi^{2}/2}\)）。
则由命题 \ref{prop:xi-bulk-curvature-exact-deconvolution-operator} 得到线性恢复算子
\[
\widehat{\nu_\sigma}(\xi)=e^{-\sigma^{2}\xi^{2}/2}\frac{\sinh(\pi\xi)}{\pi\xi}\,\widehat p(\xi).
\]
其噪声放大因子满足
\[
\mathcal A_\sigma(\xi):=e^{-\sigma^{2}\xi^{2}/2}\frac{\sinh(\pi|\xi|)}{\pi|\xi|}
\asymp
\exp\!\Bigl(\pi|\xi|-\frac{\sigma^{2}\xi^{2}}{2}\Bigr),
\]
并在 \(|\xi|\approx \pi/\sigma^{2}\) 处达到最大值 \(\asymp \exp(\pi^{2}/(2\sigma^{2}))\)。
因此在最坏情形 Fourier 噪声预算 \(\varepsilon\) 下，若希望输出噪声保持在常数级（同量级预算），则必须选取
\[
\boxed{\;\sigma\ \gtrsim\ \frac{\pi}{\sqrt{2\log(1/\varepsilon)}}.\;}
\]
\end{corollary}
\begin{proof}
由乘子表示可知输出噪声上界为 \(\sup_\xi \mathcal A_\sigma(\xi)\cdot\varepsilon\)。
指数项的最大值由一元函数 \(\pi|\xi|-\sigma^{2}\xi^{2}/2\) 的极大点 \(|\xi|=\pi/\sigma^{2}\) 给出，对应指数增益 \(\exp(\pi^{2}/(2\sigma^{2}))\)；
令 \(\varepsilon\exp(\pi^{2}/(2\sigma^{2}))\lesssim 1\) 得所示门槛。证毕。
\end{proof}

\begin{corollary}[比特复杂度下界：分辨率与精度位数的线性阈值]\label{cor:xi-bulk-curvature-bit-complexity-lower-bound}
设观测到的曲率概率密度为 \(\widetilde p=p+\eta\)，并假定其在频域带宽 \(|\xi|\le \Omega\) 上具有 \(B\) 比特精度，即
\(
\sup_{|\xi|\le\Omega}|\widehat\eta(\xi)|\le 2^{-B}
\)。
令 \(\widetilde\nu:=\mathcal D\widetilde p\)。
则由命题 \ref{prop:xi-bulk-curvature-exact-deconvolution-operator} 的乘子增长可知，在最坏情形噪声下
\[
\sup_{|\xi|\le\Omega}\bigl|\widehat{(\widetilde\nu-\nu)}(\xi)\bigr|
\ \gtrsim\
2^{-B}\,\frac{\sinh(\pi\Omega)}{\pi\Omega}
\ \asymp\
2^{-B}e^{\pi\Omega}.
\]
若要实现尺度分辨率 \(\Delta\)（通常需要 \(\Omega\gtrsim 1/\Delta\)），则必要的精度位数在量级上满足
\[
\boxed{\;B\ \gtrsim\ \frac{\pi}{\Delta\ln 2}\;}
\qquad(\text{忽略对数次级项}).
\]
\end{corollary}
\begin{proof}[证明要点]
对噪声约束选择相位对齐的 \(\widehat\eta\) 可达到最坏情形放大因子 \(\sinh(\pi\Omega)/(\pi\Omega)\)。
再以 \(\Omega\gtrsim 1/\Delta\) 换算即可得到位数阈值。证毕。
\end{proof}

\begin{remark}[屏障常数的同阶可达性：粗反卷积与 Prony 重构给出一条可执行链]\label{rem:xi-bulk-curvature-superresolution-achievability}
在有限原子口径下，一旦噪声预算满足推论 \ref{cor:xi-bulk-curvature-bit-complexity-lower-bound} 的线性位数阈值，
即可将反演实现为一条显式的可执行链：先在带宽 \(|\xi|\le \Omega\) 内施加反卷积乘子（命题 \ref{prop:xi-bulk-curvature-exact-deconvolution-operator} 或其有限阶截断版本，结论 \ref{con:xi-terminal-scale-kernel-partial-deconvolution-resolution}），
得到 \(\widehat\nu\) 的带限近似；
随后在一般位置与最小分离假设下，可用标准 Prony/矩阵铅笔/ESPRIT 等低秩谱反演由有限采样稳定恢复原子位置（与结论 \ref{con:xi-terminal-atomic-defect-hankel-prony} 的 Hankel--Prony 口径同型）。
因此常数 \(\pi/(\ln 2)\) 所给出的线性斜率不仅给出必要屏障，也在同阶上可由构造链达到（差异仅出现在对数次级项中）。
\end{remark}

\begin{remark}[\texorpdfstring{$(q=4)$}{(q=4)} 的统一阈值：涨落、反卷积与二体体积的首个可见校正阶]\label{rem:xi-scale-q4-threshold-principle}
注 \ref{rem:xi-bulk-entropy-fermi-grand-canonical} 将 bulk 熵曲率写为 Bernoulli 方差和；
在以二阶累积量（高斯近似）为基底的制度中，其首个系统性偏离由四阶累积量支配。
另一方面，结论 \ref{con:xi-terminal-scale-kernel-gamma-closed-form} 的小频率展开显示反卷积乘子 \(\log M_N(\xi)\) 的首修正项由 \(\zeta(4)\) 给出，
从而在带宽受控的粗反卷积制度中，四阶亦为最小非线性校正阶。
再者，logistic 平移族的散度字典满足 \(D_{\mathrm{KL}}=4H^2+O(\Delta^4)\)（注 \ref{rem:xi-logistic-fourth-order-closure}），
而 Hellinger--Gram 体积势的短程展开亦以 \(O(\Delta^4)\) 项给出首个超二次偏离（命题 \ref{prop:xi-logistic-hellinger-gram-determinant}）。
因此在本文的尺度全息接口中，四阶构成“首次不可忽略结构偏离”的共同闭合层级。
\end{remark}

% (User input window)
%
% 本段将 logistic 尺度核的平移族组织为一组可直接嵌入层析与判别论的闭式字典：
% 首先给出若干核心 \(f\)-散度的完全显式表达及其局部常数链；
% 随后把该核诱导的多类别 MAP 判决、最优布局、互信息与 Fisher 亏损写成可审计公式。

在命题 \ref{prop:xi-bulk-curvature-logistic-kernel-decomposition} 的尺度曲率接口中，深度原子以 logistic 核的平移族出现。
因此，平移族 \(\{\phi_a\}_{a\in\RR}\) 的散度、判决误差与信息几何常数可被视为“尺度可见性阈值”的原型账本。

\begin{proposition}[尺度曲率核的散度字典：\texorpdfstring{$1/4$}{1/4}、\texorpdfstring{$4$}{4}、\texorpdfstring{$1/3$}{1/3} 的同源封口]\label{prop:xi-logistic-divergence-dictionary}
\leanverified{paper\_xi\_logistic\_divergence\_dictionary}
令
\[
\phi(s)=\frac{1}{4\cosh^{2}(s/2)}=\frac{e^{s}}{(1+e^{s})^{2}},
\qquad
\phi_a(s):=\phi(s-a),
\qquad
\Delta:=b-a.
\]
则对任意 \(a,b\in\RR\) 下列散度均存在完全闭式：
\begin{enumerate}
  \item \textup{（总变差距离）}
  \[
  \|\phi_a-\phi_b\|_{\mathrm{TV}}
  :=\frac12\int_{\RR}|\phi_a(s)-\phi_b(s)|\,\dd s
  =\tanh\!\Bigl(\frac{|\Delta|}{4}\Bigr).
  \]
  \item \textup{（二元等先验 Bayes 最小错误率）}
  \[
  P_{\mathrm e}(\Delta)
  =\frac{1-\|\phi_a-\phi_b\|_{\mathrm{TV}}}{2}
  =\frac{1}{1+e^{|\Delta|/2}}.
  \]
  \item \textup{（Hellinger 亲和度与平方 Hellinger 距离）}
  \[
  \int_{\RR}\sqrt{\phi_a(s)\phi_b(s)}\,\dd s
  =\frac{|\Delta|}{2\sinh(|\Delta|/2)},
  \qquad
  H^{2}(\phi_a,\phi_b)
  :=1-\int_{\RR}\sqrt{\phi_a\phi_b}
  =1-\frac{|\Delta|}{2\sinh(|\Delta|/2)}.
  \]
  \item \textup{（Kullback--Leibler 散度的严格对称性）}
  \[
  D_{\mathrm{KL}}(\phi_a\|\phi_b)
  :=\int_{\RR}\phi_a(s)\log\frac{\phi_a(s)}{\phi_b(s)}\,\dd s
  =|\Delta|\coth(|\Delta|/2)-2,
  \qquad
  D_{\mathrm{KL}}(\phi_a\|\phi_b)=D_{\mathrm{KL}}(\phi_b\|\phi_a).
  \]
  \item \textup{（\(\chi^2\) 散度与 Rényi-\(2\) 散度）}
  \[
  \chi^2(\phi_a\|\phi_b)
  :=\int_{\RR}\frac{(\phi_a-\phi_b)^2}{\phi_b}\,\dd s
  =\frac{4}{3}\sinh^{2}(|\Delta|/2),
  \qquad
  D_2(\phi_a\|\phi_b)
  :=\log\!\int_{\RR}\frac{\phi_a(s)^2}{\phi_b(s)}\,\dd s
  =\log\!\Bigl(\frac{2\cosh|\Delta|+1}{3}\Bigr).
  \]
\end{enumerate}
并且在 \(|\Delta|\downarrow 0\) 时存在统一的局部量子化展开
\[
\|\phi_a-\phi_b\|_{\mathrm{TV}}=\frac{|\Delta|}{4}+O(|\Delta|^{3}),
\qquad
H^{2}(\phi_a,\phi_b)=\frac{\Delta^{2}}{24}+O(\Delta^{4}),
\qquad
D_{\mathrm{KL}}(\phi_a\|\phi_b)=\frac{\Delta^{2}}{6}+O(\Delta^{4}),
\]
从而出现严格的局部比例律
\[
D_{\mathrm{KL}}(\phi_a\|\phi_b)=4\,H^{2}(\phi_a,\phi_b)+O(\Delta^{4}),
\qquad
\|\phi_a-\phi_b\|_{\mathrm{TV}}\sim\frac{|\Delta|}{4}.
\]
该常数链把后续尺度可见性阈值固定到同一组不可消去的普适常数上。
\end{proposition}
\begin{proof}
由平移不变性，只需取 \(a=0\)、\(b=\Delta>0\)。
记 \(F(s)=\int_{-\infty}^{s}\phi(u)\dd u=(1+e^{-s})^{-1}\) 为 logistic 分布函数。
因为 \(\phi\) 为偶函数且关于 \(0\) 单峰，\(\phi(s)\) 与 \(\phi(s-\Delta)\) 的唯一交点为 \(s=\Delta/2\)，故
\[
\|\phi-\phi_{\Delta}\|_{\mathrm{TV}}
=\int_{\Delta/2}^{\infty}\bigl(\phi(s-\Delta)-\phi(s)\bigr)\dd s
=F(\Delta/2)-F(-\Delta/2)
=\tanh(\Delta/4),
\]
并由二元 Bayes 公式得到 \(P_{\mathrm e}(\Delta)\)。
\par
对 Hellinger 亲和度，写 \(\phi(s)=e^{s}(1+e^{s})^{-2}\) 并作代换 \(x=e^{s}\)（\(\dd s=\dd x/x\)）得
\[
\int_{\RR}\sqrt{\phi(s)\phi(s-\Delta)}\,\dd s
=e^{-\Delta/2}\int_{0}^{\infty}\frac{1}{(1+x)(1+e^{-\Delta}x)}\,\dd x
=\frac{\Delta}{2\sinh(\Delta/2)}.
\]
对 \(D_{\mathrm{KL}}\)，同样代换得到
\[
D_{\mathrm{KL}}(\phi\|\phi_{\Delta})
=\int_{0}^{\infty}\frac{\Delta+2\log(1+e^{-\Delta}x)-2\log(1+x)}{(1+x)^{2}}\,\dd x.
\]
对
\(
I(\alpha):=\int_{0}^{\infty}\frac{\log(1+\alpha x)}{(1+x)^{2}}\,\dd x
\)
作分部积分可得
\(
I(\alpha)=\alpha\int_{0}^{\infty}\frac{1}{(1+x)(1+\alpha x)}\dd x
=\frac{\alpha\log(1/\alpha)}{1-\alpha}
\)；
取 \(\alpha=e^{-\Delta}\) 并用 \(I(1)=1\) 得
\[
D_{\mathrm{KL}}(\phi\|\phi_{\Delta})
=\Delta+2(I(e^{-\Delta})-I(1))
=\Delta\coth(\Delta/2)-2,
\]
从而对称性由 \(\Delta\mapsto-\Delta\) 的偶性得到。
\par
对 \(\chi^2\) 与 \(D_2\)，由
\(
\chi^2(\phi\|\phi_{\Delta})=\int \phi^2/\phi_{\Delta}-1
\)，同样代换得到
\[
\int_{\RR}\frac{\phi(s)^2}{\phi(s-\Delta)}\,\dd s
=e^{\Delta}\int_{0}^{\infty}\frac{(1+e^{-\Delta}x)^{2}}{(1+x)^{4}}\,\dd x
=\frac{e^{\Delta}+1+e^{-\Delta}}{3}
=\frac{2\cosh\Delta+1}{3},
\]
其中最后一步使用 Beta 积分
\(
\int_{0}^{\infty}\frac{x^{m}}{(1+x)^{4}}\dd x=B(m+1,3-m)
\)
（\(m=0,1,2\)）。
因此
\(
\chi^2(\phi\|\phi_{\Delta})=\frac{2}{3}(\cosh\Delta-1)=\frac{4}{3}\sinh^{2}(\Delta/2)
\)
且 \(D_2=\log(1+\chi^2)\)。
局部展开由 \(\tanh\)、\(\sinh\)、\(\coth\) 的 Taylor 展开直接得到。证毕。
\end{proof}

% (User input window)

\begin{proposition}[Hellinger--Gram 体积势：平方根嵌入的闭式核与碰撞发散律]\label{prop:xi-logistic-hellinger-gram-determinant}
沿用命题 \ref{prop:xi-logistic-divergence-dictionary} 的记号并取深度谱 \(\sigma_1,\dots,\sigma_\kappa\in\RR\) 两两不同。
令
\[
\psi_j(s):=\sqrt{\phi(s-\sigma_j)}\in L^2(\RR),
\qquad
G^{(1/2)}_{jk}:=\langle \psi_j,\psi_k\rangle=\int_{\RR}\sqrt{\phi(s-\sigma_j)\phi(s-\sigma_k)}\,\dd s.
\]
则 \(G^{(1/2)}\) 正定，并且其核具有完全闭式
\[
\boxed{\;
G^{(1/2)}_{jk}=\frac{|\sigma_j-\sigma_k|}{2\sinh\!\bigl(|\sigma_j-\sigma_k|/2\bigr)}\;}
\qquad(1\le j,k\le\kappa),
\]
其中对角线取连续延拓值 \(G^{(1/2)}_{jj}=1\)。
定义体积与熵势
\[
V_{1/2}:=\sqrt{\det G^{(1/2)}},
\qquad
S_{1/2}:=-\log\det G^{(1/2)}.
\]
当发生尺度碰撞 \(\sigma_j=\sigma_\star+\varepsilon t_j\) 且 \(\varepsilon\downarrow 0\)（\(t_j\) 两两不同）时，有普适 Vandermonde 发散律
\[
\boxed{\;
S_{1/2}(\boldsymbol{\sigma})
=\kappa(\kappa-1)\log\frac{1}{\varepsilon}
+2\sum_{1\le j<k\le \kappa}\log\frac{1}{|t_k-t_j|}
-\log\Bigl(\mathfrak H_\kappa\prod_{j=0}^{\kappa-1}\frac{1}{(j!)^{2}}\Bigr)
+o(1).\;}
\]
其中 \(\mathfrak H_\kappa\) 为定理 \ref{thm:xi-hellinger-gram-collision-hankel-constant} 所定义的普适 Hankel 常数（仅依赖于 \(\kappa\) 与谱密度 \(w(\xi)=\pi^2\mathrm{sech}^2(\pi\xi)\)）。
其中二体对数项系数恒为 \(2\)，与谱点权重及其余构型细节无关。
\end{proposition}
\begin{proof}
正定性由 \(G^{(1/2)}\) 的 Gram 表示直接得到。
闭式由命题 \ref{prop:xi-logistic-divergence-dictionary} 第 \textup{(3)} 条（Hellinger 亲和度闭式）取 \(a=\sigma_j\)、\(b=\sigma_k\) 立即推出。
\par
对碰撞渐近，定理 \ref{thm:xi-hellinger-gram-collision-hankel-constant} 给出更精确的行列式渐近：
\[
\det G^{(1/2)}
=\varepsilon^{\kappa(\kappa-1)}
\Bigl(\prod_{j=0}^{\kappa-1}\frac{1}{(j!)^{2}}\Bigr)\,
\mathfrak H_\kappa\,
\prod_{1\le j<k\le\kappa}(t_k-t_j)^{2}\,\bigl(1+o(1)\bigr).
\]
取负对数即得所示展开。证毕。
\end{proof}

% (User input window)
%
% 本段将 Hellinger 二体核 \(A(\Delta)\) 由闭式亲和度提升为一条完全谱确定的几何对象：
% 其 Fourier 密度为 \(\cosh^{-2}(\pi\xi)\)，并诱导一个显式 RKHS 与行列式体积势；
% 在等距谱下进一步 Toeplitz 化并给出 Szeg\H{o} 极限与 CUE 表示。

\begin{proposition}[Hellinger 平移核的 Fourier 完全闭式：\texorpdfstring{$\cosh^{-2}$}{cosh^{-2}} 为深度二体几何的谱密度]\label{prop:xi-hellinger-kernel-fourier-sech2}
\leanverified{paper\_xi\_hellinger\_kernel\_fourier\_sech2}
令尺度单峰核
\[
\phi(s)=\frac{1}{4\cosh^{2}(s/2)}\qquad(s\in\RR),
\]
并定义 Hellinger 亲和核
\[
A(\Delta):=\int_{\RR}\sqrt{\phi(s)\phi(s-\Delta)}\,\dd s
=\frac{|\Delta|}{2\sinh(|\Delta|/2)}\qquad(\Delta\in\RR).
\]
则 \(A\in L^{1}(\RR)\cap C^\infty(\RR)\) 为偶函数。
在 Fourier 变换约定
\(
\widehat f(\xi):=\int_{\RR}e^{-\mathrm{i}\xi x}f(x)\,\dd x
\)
下，有严格恒等式
\[
\boxed{\;
\widehat A(\xi)=\int_{\RR}A(\Delta)e^{-\mathrm{i}\xi\Delta}\,\dd\Delta
=\frac{\pi^{2}}{\cosh^{2}(\pi\xi)}\;}
\qquad(\xi\in\RR).
\]
特别地
\(
\int_{\RR}A(\Delta)\,\dd\Delta=\widehat A(0)=\pi^{2}
\)。
\end{proposition}
\begin{proof}
令
\(
f(s):=\sqrt{\phi(s)}=\frac{1}{2\cosh(s/2)}
\)。
则
\(
A(\Delta)=\int_{\RR}f(s)f(s-\Delta)\,\dd s
\)
为 \(f\) 的自相关，因而 \(\widehat A(\xi)=|\widehat f(\xi)|^{2}\)。
作代换 \(s=2u\) 得
\[
\widehat f(\xi)
=\int_{\RR}\frac{1}{2\cosh(s/2)}e^{-\mathrm{i}\xi s}\,\dd s
=\int_{\RR}\frac{1}{\cosh u}\,e^{-2\mathrm{i}\xi u}\,\dd u
=\frac{\pi}{\cosh(\pi\xi)},
\]
其中最后一步为 \(\cosh^{-1}\) 的 Fourier 闭式。
故 \(\widehat A(\xi)=\pi^{2}/\cosh^{2}(\pi\xi)\)。证毕。
\end{proof}

\input{para__xi-hellinger-green-operator-energy-duality}

\begin{proposition}[深度谱的 RKHS 化：\texorpdfstring{$A$}{A} 为严格再生核且 \texorpdfstring{$\det G$}{det G} 为评价函子体积]\label{prop:xi-hellinger-kernel-rkhs-volume}
沿用命题 \ref{prop:xi-hellinger-kernel-fourier-sech2} 的记号，并记谱密度
\[
w(\xi):=\widehat A(\xi)=\frac{\pi^2}{\cosh^2(\pi\xi)}.
\]
定义 Hilbert 空间
\[
\mathcal H_{A}
:=\Bigl\{f\in\mathcal S'(\RR):\ \|f\|_{\mathcal H_A}^{2}
:=\frac{1}{2\pi}\int_{\RR}\frac{|\widehat f(\xi)|^{2}}{w(\xi)}\,\dd\xi
=\frac{1}{2\pi}\int_{\RR}|\widehat f(\xi)|^{2}\frac{\cosh^{2}(\pi\xi)}{\pi^{2}}\,\dd\xi<\infty\Bigr\}.
\]
则对任意 \(x\in\RR\)，评价泛函 \(f\mapsto f(x)\) 在 \(\mathcal H_A\) 上连续。
令 \(k_x(\cdot):=A(\cdot-x)\)，则 \(k_x\in\mathcal H_A\) 且对所有 \(f\in\mathcal H_A\) 有再生性质
\[
\boxed{\;f(x)=\langle f,k_x\rangle_{\mathcal H_A}\;},
\qquad
\langle k_x,k_y\rangle_{\mathcal H_A}=A(x-y).
\]
给定深度对数谱 \(\boldsymbol{\sigma}=(\sigma_1,\dots,\sigma_\kappa)\)，定义 Gram 矩阵
\[
G(\boldsymbol{\sigma})_{jk}:=A(\sigma_j-\sigma_k)\qquad(1\le j,k\le\kappa).
\]
则
\[
\boxed{\;
\det G(\boldsymbol{\sigma})
=\mathrm{Vol}^{2}\bigl(\mathrm{span}\{k_{\sigma_1},\dots,k_{\sigma_\kappa}\}\bigr),\qquad
S_{1/2}(\boldsymbol{\sigma}):=-\log\det G(\boldsymbol{\sigma}).\;}
\]
\end{proposition}
\begin{proof}
在 \(\mathcal H_A\) 上定义内积
\(
\langle f,g\rangle_{\mathcal H_A}:=\frac{1}{2\pi}\int \widehat f(\xi)\overline{\widehat g(\xi)}\,w(\xi)^{-1}\dd\xi
\)。
对 \(k_x(\cdot)=A(\cdot-x)\)，由 Fourier 平移得 \(\widehat{k_x}(\xi)=w(\xi)e^{-\mathrm{i}\xi x}\)，从而 \(k_x\in\mathcal H_A\) 且
\[
\langle f,k_x\rangle_{\mathcal H_A}
=\frac{1}{2\pi}\int_{\RR}\widehat f(\xi)\,e^{\mathrm{i}\xi x}\,\dd\xi
=f(x),
\]
其中最后一步为 Schwartz 分布的 Fourier 反演公式。
取 \(f=k_x\) 得 \(\langle k_x,k_y\rangle_{\mathcal H_A}=A(x-y)\)。
最后一式为 Hilbert 空间中 Gram 行列式的几何解释：\(\det(\langle v_j,v_k\rangle)=\mathrm{Vol}^{2}(\mathrm{span}\{v_j\})\)。证毕。
\end{proof}

\input{para__xi-gram-volume-monotonicity-coherence-nearestneighbor}

\begin{proposition}[深度体积熵势的精确一阶变分与平衡方程]\label{prop:xi-hellinger-gram-first-variation-equilibrium}
令深度对数谱 \(\mathbf s=(s_1,\dots,s_\kappa)\) 两两不同，并令 Hellinger--Gram 矩阵
\[
G(\mathbf s)_{jk}=A(s_j-s_k)\qquad(1\le j,k\le \kappa),
\]
其中核 \(A\) 见命题 \ref{prop:xi-hellinger-kernel-fourier-sech2}。
定义深度体积熵势
\[
S_{1/2}(\mathbf s):=-\log\det G(\mathbf s).
\]
则对每个 \(1\le j\le \kappa\) 有严格梯度公式
\[
\boxed{\;
\partial_{s_j}S_{1/2}(\mathbf s)
=-2\sum_{k=1}^{\kappa}\bigl(G(\mathbf s)^{-1}\bigr)_{jk}\,A'(s_j-s_k)\;}
\]
其中 \(A'(0):=0\)，且对 \(\Delta\neq 0\),
\[
\boxed{\;
A'(\Delta)
=A(\Delta)\Bigl(\frac{1}{\Delta}-\frac12\coth(\Delta/2)\Bigr)\; }.
\]
由于 \(S_{1/2}\) 对整体平移 \(\mathbf s\mapsto \mathbf s+c\mathbf 1\) 不变，任一在平移自由度被固定后的驻点必须满足平衡方程
\[
\sum_{k=1}^{\kappa}\bigl(G(\mathbf s)^{-1}\bigr)_{jk}\,A'(s_j-s_k)=0,\qquad j=1,\dots,\kappa.
\]
并且力核 \(\frac{1}{\Delta}-\frac12\coth(\Delta/2)\) 在 \(|\Delta|\to\infty\) 与 \(\Delta\to 0\) 的两端满足
\[
\frac{1}{\Delta}-\frac12\coth(\Delta/2)= -\frac12\,\mathrm{sgn}(\Delta)+O(|\Delta|^{-1}),
\qquad
\frac{1}{\Delta}-\frac12\coth(\Delta/2)=-\frac{\Delta}{12}+O(\Delta^3).
\]
\end{proposition}
\begin{proof}
由矩阵微分恒等式 \(\partial\log\det G=\mathrm{tr}(G^{-1}\partial G)\)，得
\[
\partial_{s_j}\log\det G(\mathbf s)=\mathrm{tr}\bigl(G(\mathbf s)^{-1}\,\partial_{s_j}G(\mathbf s)\bigr).
\]
注意 \(\partial_{s_j}G_{ab}=\delta_{aj}A'(s_j-s_b)-\delta_{bj}A'(s_a-s_j)\)，且 \(A'\) 为奇函数。
于是
\[
\partial_{s_j}\log\det G(\mathbf s)
=\sum_{k=1}^{\kappa}\bigl(G^{-1}\bigr)_{jk}A'(s_j-s_k)
+\sum_{k=1}^{\kappa}\bigl(G^{-1}\bigr)_{kj}A'(s_j-s_k)
=2\sum_{k=1}^{\kappa}\bigl(G^{-1}\bigr)_{jk}A'(s_j-s_k),
\]
从而得到所示 \(\partial_{s_j}S_{1/2}=-\partial_{s_j}\log\det G\) 公式。
对 \(\Delta\neq 0\)，直接对 \(A(\Delta)=\Delta/(2\sinh(\Delta/2))\) 求导并整理得
\(
A'(\Delta)=A(\Delta)\bigl(\Delta^{-1}-\tfrac12\coth(\Delta/2)\bigr)
\)；再由奇对称令 \(A'(0)=0\)。
\par
平衡方程来自驻点条件与平移不变性（等价于在约束 \(\sum_j s_j=0\) 或固定任一基点的条件下取驻点）。
两端展开由 \(\coth x=1+O(e^{-2x})\)（\(x\to\infty\)）与
\(\coth x=x^{-1}+x/3+O(x^3)\)（\(x\to 0\)）代入即得。证毕。
\end{proof}

\begin{remark}[两体势的面积律与对数修正：\texorpdfstring{$1/4$}{1/4} 的长程再现]\label{rem:xi-hellinger-two-body-potential-area-log}
令两体势 \(V(\Delta):=-\log A(\Delta)\)（\(\Delta\ge 0\)）。
则对 \(\Delta>0\) 有严格恒等式
\[
\boxed{\;
V(\Delta)=\frac{\Delta}{2}-\log\Delta+\log\bigl(1-e^{-\Delta}\bigr)\;}
\]
并因此
\[
V(\Delta)=\frac{\Delta}{2}-\log\Delta+O(e^{-\Delta})\qquad(\Delta\to\infty).
\]
若改用物理距离变量 \(d:=2\Delta\)，则
\[
V(d/2)=\frac{d}{4}-\log d+\log 2+O(e^{-d/2})\qquad(d\to\infty),
\]
其中线性项系数 \(1/4\) 在该核的长程势中以严格形式出现，而对数项给出不可消去的次级校正。
\end{remark}

\begin{proposition}[深度能量泛函的严格凸性：\texorpdfstring{$-\log A$}{-log A} 的条件负定性]\label{prop:xi-hellinger-energy-functional-strict-convexity}
定义连续能量泛函
\[
\mathcal E(\nu):=\frac12\iint_{\RR\times\RR} -\log A(x-y)\,\dd\nu(x)\dd\nu(y),
\]
其中 \(\nu\) 为具有有限能量的概率测度。
则核 \(\psi(\Delta):=-\log A(\Delta)\) 为（严格）条件负定函数：对任意实系数 \(c_1,\dots,c_m\) 满足 \(\sum_i c_i=0\) 与任意点 \(x_1,\dots,x_m\in\RR\) 有
\[
\sum_{i,j=1}^{m}c_ic_j\,\psi(x_i-x_j)\le 0,
\]
并因此 \(\mathcal E\) 在概率测度的凸集合上为（严格）凸泛函。
在消除整体平移自由度的任一仿射约束下（例如固定均值），\(\mathcal E\) 的极小测度唯一。
\end{proposition}
\begin{proof}[证明要点]
由 Weierstrass 乘积
\(
2\sinh(\Delta/2)=\Delta\prod_{n\ge 1}\bigl(1+\Delta^2/(4\pi^2n^2)\bigr)
\)
得
\[
A(\Delta)=\prod_{n\ge 1}\Bigl(1+\frac{\Delta^2}{4\pi^2n^2}\Bigr)^{-1}.
\]
每个因子 \((1+b^2\Delta^2)^{-1}\) 是中心 Laplace 分布的特征函数，故 \(A\) 为一族独立 Laplace 变量之和的特征函数，从而为无穷可分；由 L\'evy--Khintchine 理论，\(\psi=-\log A\) 为连续负定函数（等价于所述条件负定不等式）。
\par
对 \(\nu_t=(1-t)\nu_0+t\nu_1\) 展开 \(\mathcal E(\nu_t)\) 可得
\(
\mathcal E(\nu_t)=(1-t)\mathcal E(\nu_0)+t\mathcal E(\nu_1)-\frac{t(1-t)}{2}\iint \psi\,\dd(\nu_1-\nu_0)\dd(\nu_1-\nu_0)
\)；
条件负定性使最后一项非负，从而得到凸性。严格性与约束下的唯一性由等号情形强制 \(\nu_0=\nu_1\)（或仅差整体平移）推出。证毕。
\end{proof}

\begin{theorem}[深度 Hellinger--Gram 的严格全正性与谱振荡律]\label{thm:xi-hellinger-gram-strict-total-positivity}
\leanverified{paper\_xi\_hellinger\_gram\_strict\_total\_positivity}
对严格递增深度 \(s_1<\cdots<s_\kappa\)，矩阵 \(G(\mathbf s)=(A(s_j-s_k))_{1\le j,k\le\kappa}\) 为严格全正矩阵。
因此：
\begin{itemize}
\item 全部特征值均为单根；
\item 第 \(r\) 大特征值对应的特征向量在索引序上恰有 \(r-1\) 次严格符号变号；
\item 任意非零向量经 \(G(\mathbf s)\) 作用后的符号变号数严格不增（变差递减）。
\end{itemize}
\end{theorem}
\begin{proof}[证明要点]
由命题 \ref{prop:xi-hellinger-kernel-fourier-sech2}，归一化特征函数满足
\[
\frac{\widehat A(\xi)}{\widehat A(0)}=\mathrm{sech}^{2}(\pi\xi)
=\prod_{n\ge 1}\Bigl(1+\frac{\xi^{2}}{(n-\tfrac12)^{2}}\Bigr)^{-2},
\]
其中使用 \(\cosh(\pi z)=\prod_{n\ge 1}(1+z^{2}/(n-\tfrac12)^{2})\) 的 Weierstrass 乘积。
因此 \(A/\widehat A(0)\) 为 PF\(_\infty\)（P\'olya frequency）密度；由 Schoenberg 的 PF\(_\infty\)\(\Rightarrow\)严格全正平移核定理可知，核 \((x,y)\mapsto A(x-y)\) 为严格全正。
令 \((x_i)=(y_i)=(s_i)\) 即得 \(G(\mathbf s)\) 的严格全正性。
谱单纯性、特征向量振荡律与变差递减性质为全正矩阵的基本定理（Gantmacher--Kre\u{\i}n 变差递减/振荡定理）的直接推论。证毕。
\end{proof}

\begin{corollary}[逆矩阵的棋盘符号律：交替耦合的代数证书]\label{cor:xi-hellinger-gram-inverse-checkerboard}
在定理 \ref{thm:xi-hellinger-gram-strict-total-positivity} 的条件下，\(G(\mathbf s)^{-1}\) 具有严格棋盘符号：
\[
\boxed{\;
(-1)^{i+j}\bigl(G(\mathbf s)^{-1}\bigr)_{ij}>0\qquad(1\le i,j\le \kappa).\;}
\]
\end{corollary}
\begin{proof}
这是全正矩阵逆的符号正规性（sign-regular）结论：非奇异全正矩阵的逆为严格棋盘符号矩阵。证毕。
\end{proof}

\begin{proposition}[Andreief 表示与频域对数气体：\texorpdfstring{$S_{1/2}$}{S_{1/2}} 为外势自由能]\label{prop:xi-hellinger-kernel-andreief-loggas}
在命题 \ref{prop:xi-hellinger-kernel-rkhs-volume} 的设定下，有严格积分表示
\[
\boxed{\;
\det G(\boldsymbol{\sigma})
=\frac{1}{\kappa!(2\pi)^{\kappa}}
\int_{\RR^\kappa}
\Bigl(\prod_{r=1}^{\kappa} w(\xi_r)\Bigr)\,
\bigl|\det(e^{-\mathrm{i}\xi_r\sigma_j})_{1\le r,j\le\kappa}\bigr|^{2}
\ \dd\xi_1\cdots \dd\xi_\kappa.\;}
\]
因此 \(S_{1/2}(\boldsymbol{\sigma})\) 等价于一维 \(\beta=2\) 对数相互作用体系在外势
\(
V(\xi):=-\log w(\xi)=2\log\cosh(\pi\xi)-\log\pi^{2}
\)
下的配分函数自由能；深度谱 \(\boldsymbol{\sigma}\) 以外部相位插值进入 Vandermonde 型体积因子。
\end{proposition}
\begin{proof}
由命题 \ref{prop:xi-hellinger-kernel-fourier-sech2} 的 Fourier 表示，
\[
A(x-y)=\frac{1}{2\pi}\int_{\RR}w(\xi)\,e^{\mathrm{i}\xi(x-y)}\,\dd\xi.
\]
令测度 \(\dd\mu(\xi):=w(\xi)\dd\xi/(2\pi)\)，并置
\(
F_j(\xi):=e^{-\mathrm{i}\xi\sigma_j}
\)。
则 \(G_{jk}=\int F_j(\xi)\overline{F_k(\xi)}\,\dd\mu(\xi)\)。
由 Andreief 恒等式（亦即 Gram 行列式的积分化）得到
\[
\det(G_{jk})
=\frac{1}{\kappa!}\int_{\RR^\kappa}\bigl|\det(F_j(\xi_r))_{1\le r,j\le\kappa}\bigr|^{2}\,\dd\mu(\xi_1)\cdots\dd\mu(\xi_\kappa),
\]
化为所示形式。自由能表述由对数取负得到。证毕。
\end{proof}

\begin{theorem}[碰撞常数的 Hankel 封口：簇化缩放下 \texorpdfstring{$\det G$}{det G} 的 confluent Vandermonde 因子化]\label{thm:xi-hellinger-gram-collision-hankel-constant}
沿用命题 \ref{prop:xi-hellinger-kernel-andreief-loggas} 的记号，并取任意 \(\kappa\ge 1\)。
给定两两不同的 \(t_1,\dots,t_\kappa\in\RR\) 与任意 \(s_0\in\RR\)，考虑簇化缩放
\[
\sigma_j=s_0+\varepsilon t_j\qquad(1\le j\le \kappa),\qquad \varepsilon\downarrow 0,
\]
并令 \(G(\boldsymbol{\sigma})=(A(\sigma_j-\sigma_k))_{1\le j,k\le\kappa}\)。
则存在严格渐近
\[
\boxed{\;
\det G(\boldsymbol{\sigma})
=\varepsilon^{\kappa(\kappa-1)}
\Bigl(\prod_{j=0}^{\kappa-1}\frac{1}{(j!)^{2}}\Bigr)\,
\mathfrak H_\kappa\,
\prod_{1\le j<k\le\kappa}(t_k-t_j)^{2}\,\bigl(1+o(1)\bigr)\;}
\]
其中普适常数 \(\mathfrak H_\kappa\) 为权函数 \(w(\xi)=\pi^{2}\mathrm{sech}^{2}(\pi\xi)\) 的正规化 Hankel 行列式：
\[
\mathfrak H_\kappa:=\det(\mu_{p+q})_{p,q=0}^{\kappa-1},\qquad
\mu_n:=\frac{1}{2\pi}\int_{\RR}\xi^{n}w(\xi)\,\dd\xi.
\]
因此碰撞发散的 Vandermonde 幂律之外，仍存在仅依赖于 \(\kappa\) 的不可约体积常数 \(\mathfrak H_\kappa\)。
\end{theorem}
\begin{proof}
由命题 \ref{prop:xi-hellinger-kernel-andreief-loggas} 的 Andreief 表示与平移不变性，\(\det G(\boldsymbol{\sigma})\) 仅依赖差分 \(\sigma_j-\sigma_k=\varepsilon(t_j-t_k)\)，且
\[
\det G(\boldsymbol{\sigma})
=\frac{1}{\kappa!(2\pi)^{\kappa}}
\int_{\RR^\kappa}\Bigl(\prod_{r=1}^{\kappa}w(\xi_r)\Bigr)\,
\bigl|\det(e^{-\mathrm{i}\varepsilon\xi_r t_j})_{1\le r,j\le\kappa}\bigr|^{2}\,\dd\xi_1\cdots\dd\xi_\kappa.
\]
当 \(\varepsilon\downarrow 0\) 时，经典的 confluent Vandermonde 渐近给出
\[
\det(e^{-\mathrm{i}\varepsilon\xi_r t_j})_{r,j}
=\varepsilon^{\kappa(\kappa-1)/2}\Bigl(\prod_{j=0}^{\kappa-1}\frac{1}{j!}\Bigr)
\Delta(\boldsymbol{\xi})\,\Delta(\mathbf t)\,\bigl(1+o(1)\bigr),
\]
其中 \(\Delta(\boldsymbol{\xi})=\prod_{1\le r<\ell\le\kappa}(\xi_\ell-\xi_r)\)、\(\Delta(\mathbf t)=\prod_{1\le j<k\le\kappa}(t_k-t_j)\)。
取模平方并代回，得到
\[
\det G(\boldsymbol{\sigma})
=\varepsilon^{\kappa(\kappa-1)}\Bigl(\prod_{j=0}^{\kappa-1}\frac{1}{(j!)^{2}}\Bigr)\,
\Delta(\mathbf t)^{2}\cdot
\frac{1}{\kappa!}\int_{\RR^\kappa}\Delta(\boldsymbol{\xi})^{2}\,\dd\mu(\xi_1)\cdots\dd\mu(\xi_\kappa)\,\bigl(1+o(1)\bigr),
\]
其中 \(\dd\mu(\xi)=w(\xi)\dd\xi/(2\pi)\)。
最后，由 Andreief 恒等式（Vandermonde 平方的矩阵化积分）有
\[
\frac{1}{\kappa!}\int_{\RR^\kappa}\Delta(\boldsymbol{\xi})^{2}\,\dd\mu(\xi_1)\cdots\dd\mu(\xi_\kappa)
=\det(\mu_{p+q})_{p,q=0}^{\kappa-1}=\mathfrak H_\kappa,
\]
从而得到所示渐近。证毕。
\end{proof}

\begin{corollary}[簇化相的最小奇异指数：条件数的制度性幂次爆炸]\label{cor:xi-hellinger-gram-cluster-conditioning-exponent}
在定理 \ref{thm:xi-hellinger-gram-collision-hankel-constant} 的簇化缩放下，
存在常数 \(c(\mathbf t)>0\) 使
\[
\det G(\boldsymbol{\sigma})=c(\mathbf t)\,\varepsilon^{\kappa(\kappa-1)}\bigl(1+o(1)\bigr).
\]
并且由于 \(\mathrm{tr}(G)=\kappa\)，最大特征值满足 \(\lambda_{\max}(G)\le \kappa\)。
因此最小特征值具有硬上界
\[
\lambda_{\min}(G)\le \frac{\det G}{\lambda_{\max}(G)^{\kappa-1}}
\le \kappa^{1-\kappa}\,c(\mathbf t)\,\varepsilon^{\kappa(\kappa-1)}\bigl(1+o(1)\bigr),
\]
从而谱条件数满足
\[
\boxed{\;
\mathrm{cond}_2(G)=\frac{\lambda_{\max}(G)}{\lambda_{\min}(G)}
\ \ge\ C(\mathbf t)\,\varepsilon^{-\kappa(\kappa-1)}\bigl(1+o(1)\bigr)\;}
\qquad(\varepsilon\downarrow 0),
\]
其中 \(C(\mathbf t)>0\) 为常数。
因此簇化相的数值不适定性由指数 \(\kappa(\kappa-1)\) 严格锁定。
\end{corollary}
\begin{proof}
行列式与谱乘积给出
\(
\det G=\lambda_{\min}\prod_{r\neq\min}\lambda_r\ge \lambda_{\min}\lambda_{\max}^{\kappa-1}
\)，
故 \(\lambda_{\min}\le \det G/\lambda_{\max}^{\kappa-1}\)。
又 \(\lambda_{\max}\le \mathrm{tr}(G)=\kappa\)，并代入定理 \ref{thm:xi-hellinger-gram-collision-hankel-constant} 的渐近，即得所示条件数下界。证毕。
\end{proof}

\begin{proposition}[普适 Hankel 常数的算术性：\texorpdfstring{$\mathfrak H_\kappa\in\QQ$}{H_k in Q} 且由偶 \texorpdfstring{$\zeta$}{zeta} 值显式封口]\label{prop:xi-hellinger-hankel-constant-rationality}
\leanverified{paper\_xi\_hellinger\_hankel\_constant\_rationality}
沿用定理 \ref{thm:xi-hellinger-gram-collision-hankel-constant} 的记号。
则矩 \(\mu_n\) 满足 \(\mu_{2n+1}=0\)，且对任意 \(n\ge 0\) 有严格闭式
\[
\boxed{\;
\mu_{2n}
=2^{1-2n}\pi^{-2n}(2n)!\,\bigl(1-2^{1-2n}\bigr)\zeta(2n)\in\QQ\; }.
\]
从而对任意 \(\kappa\ge 1\)，\(\mathfrak H_\kappa=\det(\mu_{p+q})_{p,q=0}^{\kappa-1}\in\QQ\)。
例如
\[
\boxed{\;
\mathfrak H_1=1,\quad
\mathfrak H_2=\frac{1}{12},\quad
\mathfrak H_3=\frac{1}{540},\quad
\mathfrak H_4=\frac{1}{42000},\quad
\mathfrak H_5=\frac{1}{3215625}.\;}
\]
\end{proposition}
\begin{proof}[证明要点]
奇次矩为零来自 \(w(\xi)\) 的偶对称。
对偶次矩，可由 \(\mathrm{sech}^2\) 的部分分式/留数展开（等价地，由 \(\tanh\) 的晶格极点结构）将
\(\int_{\RR}\xi^{2n}\mathrm{sech}^{2}(\pi\xi)\dd\xi\)
化为 Hurwitz \(\zeta\) 值之差，从而得到所示 \(\zeta(2n)\) 闭式；再由
\(\zeta(2n)\in \pi^{2n}\QQ\)
推出 \(\mu_{2n}\in\QQ\)。
由 Hankel 行列式对矩的多项式性即得 \(\mathfrak H_\kappa\in\QQ\)。
所列数值由该闭式直接计算。证毕。
\end{proof}

\begin{remark}[算术带宽定律：\texorpdfstring{$\mathfrak H_\kappa$}{H_k} 的 \texorpdfstring{$p$}{p}-进刚性仅支持于 \texorpdfstring{$p\le 2\kappa-1$}{p<=2k-1}]\label{rem:xi-hellinger-hankel-padic-bandwidth}
令 \(\nu_p(\cdot)\) 为 \(p\)-进赋值。
由命题 \ref{prop:xi-hellinger-hankel-constant-rationality} 中 \(\mu_{2n}\) 的 Bernoulli--von Staudt--Clausen 分母结构可知：
对任意素数 \(p>2\kappa-1\)，有 \(\nu_p(\mu_0),\dots,\nu_p(\mu_{2\kappa-2})\ge 0\)。
因此 \(\nu_p(\mathfrak H_\kappa)=0\)（\(p>2\kappa-1\)）。
换言之，\(\mathfrak H_\kappa\) 的全部 \(p\)-进“病态/刚性”只能由有限素数集合 \(\{p:\ p\le 2\kappa-1\}\) 贡献。
\end{remark}

\begin{proposition}[等距深度的 Toeplitz 化：\texorpdfstring{$\cosh^{-2}$}{cosh^{-2}} 的 Poisson 周期和符号]\label{prop:xi-hellinger-toeplitz-symbol-poisson}
若 \(\sigma_j=\sigma_0+(j-1)\Delta\)（\(\Delta>0\)），则 \(G(\boldsymbol{\sigma})\) 为 Toeplitz 矩阵
\(
G_{jk}=A((j-k)\Delta)
\)。
其 Toeplitz 符号 \(\sigma_\Delta:[-\pi,\pi]\to(0,\infty)\) 可写为
\[
\boxed{\;
\sigma_\Delta(\theta)
:=\sum_{n\in\ZZ}A(n\Delta)e^{\mathrm{i}n\theta}
=\frac{1}{\Delta}\sum_{k\in\ZZ}\frac{\pi^{2}}{\cosh^{2}\!\Bigl(\pi\frac{\theta+2\pi k}{\Delta}\Bigr)}\;}
\qquad(\theta\in[-\pi,\pi]).
\]
等价地，令 \(q:=e^{-\Delta/2}\) 并取 Jacobi \(\vartheta_4\) 函数（nome \(q\)）
\[
\vartheta_4(z,q):=\sum_{m\in\ZZ}(-1)^m q^{m^2}e^{2\mathrm{i}mz},
\]
则对所有 \(\theta\in[-\pi,\pi]\) 有严格恒等式
\[
\boxed{\;
\sigma_\Delta(\theta)
=1+\Delta\,\partial_\theta^{2}\log\vartheta_4\!\Bigl(\frac{\theta}{2},q\Bigr)
=1+2\Delta\sum_{n\ge 1}\frac{nq^{n}}{1-q^{2n}}\cos(n\theta).\;}
\]
\end{proposition}
\begin{proof}
Toeplitz 性由 \(G_{jk}=A(\sigma_j-\sigma_k)=A((j-k)\Delta)\) 直接得到。
Poisson 求和公式在本 Fourier 约定下给出
\(
\sum_{n\in\ZZ}A(n\Delta)e^{\mathrm{i}n\theta}
=\Delta^{-1}\sum_{k\in\ZZ}\widehat A((\theta+2\pi k)/\Delta)
\)。
代入命题 \ref{prop:xi-hellinger-kernel-fourier-sech2} 的 \(\widehat A=w\) 即得 Poisson 周期和表示。
\par
另一方面，由
\(
A(n\Delta)=\frac{n\Delta}{2\sinh(n\Delta/2)}
=n\Delta\,\frac{q^n}{1-q^{2n}}
\)
（\(n\ge 1\)、\(q=e^{-\Delta/2}\)）得 Fourier 级数表示
\[
\sigma_\Delta(\theta)=1+2\sum_{n\ge 1}A(n\Delta)\cos(n\theta)
=1+2\Delta\sum_{n\ge 1}\frac{nq^{n}}{1-q^{2n}}\cos(n\theta).
\]
再由 \(\vartheta_4\) 的 Jacobi 乘积公式推出标准恒等式
\[
\partial_z^{2}\log\vartheta_4(z,q)
=8\sum_{n\ge 1}\frac{nq^{n}}{1-q^{2n}}\cos(2nz),
\]
取 \(z=\theta/2\) 并用 \(\partial_\theta^{2}=(1/4)\partial_z^{2}\) 即得导数闭式。证毕。
\end{proof}

\begin{remark}[椭圆延拓与半周期二阶极点晶格]\label{rem:xi-hellinger-toeplitz-elliptic-extension-halfperiod-poles}
把 \(\sigma_\Delta(\theta)\) 视为复变量 \(\theta\) 的函数。
由 \(\vartheta_4\) 的准周期性可知 \(\partial_\theta^2\log\vartheta_4(\theta/2,q)\) 消去仿射漂移项，从而 \(\sigma_\Delta\) 具有基本周期 \(2\pi\) 与 \(\mathrm{i}\Delta\)。
记矩形晶格
\[
\Lambda_\Delta:=2\pi\ZZ+\mathrm{i}\Delta\ZZ,
\qquad
E_\Delta:=\CC/\Lambda_\Delta.
\]
则 \(\sigma_\Delta\) 在椭圆曲线 \(E_\Delta\) 上诱导一个亚纯函数；
并且由命题 \ref{prop:xi-hellinger-toeplitz-symbol-poisson} 的 Poisson 周期和表示可知其极点均为二阶极点且恰位于半周期平移晶格
\[
\theta\in \mathrm{i}\Delta\Bigl(\ZZ+\tfrac12\Bigr)+2\pi\ZZ.
\]
\end{remark}

\begin{remark}[模对偶尺度 \texorpdfstring{$\Delta\leftrightarrow 2\pi^{2}/\Delta$}{Delta <-> 2pi^2/Delta}：Poisson 尖峰与 Fourier 衰减的强制等价]\label{rem:xi-hellinger-toeplitz-modular-duality-scale}
在命题 \ref{prop:xi-hellinger-toeplitz-symbol-poisson} 中令 \(\tau:=\mathrm{i}\Delta/(2\pi)\)，则 nome \(q=e^{\mathrm{i}\pi\tau}=e^{-\Delta/2}\)。
Jacobi 变换 \(\tau\mapsto -1/\tau\) 导出对偶 nome
\[
q_\star:=e^{\mathrm{i}\pi(-1/\tau)}=e^{-2\pi^{2}/\Delta}.
\]
在该变换下，\(\sigma_\Delta\) 的 Fourier 级数表示与 Poisson 周期和表示互为同一 θ-结构的两种等价表象；
其中尺度常数 \(2\pi^{2}/\Delta\) 并非近似技巧，而由 θ-函数的模对偶性制度性地强制出现。
\end{remark}

\begin{proposition}[等距 Toeplitz 相的逆矩阵局域化：反演算子的指数衰减由带宽封口]\label{prop:xi-hellinger-toeplitz-inverse-localization}
在命题 \ref{prop:xi-hellinger-toeplitz-symbol-poisson} 的等距设定下，令
\(
G_\kappa(\Delta)=(A((j-k)\Delta))_{1\le j,k\le \kappa}
\)
并记其符号为 \(\sigma_\Delta\)。
则 \(\sigma_\Delta\) 在带状域 \(|\Im\theta|<\Delta/2\) 内解析。
并且对任意 \(0<\varepsilon<\Delta/2\)，存在常数 \(C_{\Delta,\varepsilon}>0\) 使对一切 \(\kappa\) 与 \(1\le j,k\le\kappa\) 有
\[
\boxed{\;
\bigl|\bigl(G_\kappa(\Delta)^{-1}\bigr)_{jk}\bigr|
\le C_{\Delta,\varepsilon}\,e^{-(\Delta/2-\varepsilon)|j-k|}\;}
\]
从而等距深度编码具有严格的“局域可逆性”。
指数门槛 \(\Delta/2\) 由 \(\sigma_\Delta\) 的最近复奇点距离封口（见命题 \ref{prop:xi-hellinger-toeplitz-symbol-poisson} 的 Fourier 系数衰减率）。
\end{proposition}
\begin{proof}[证明要点]
由 \(A(n\Delta)=O(n\Delta\,e^{-n\Delta/2})\) 得 \(\sigma_\Delta\) 的 Fourier 系数指数衰减，从而 \(\sigma_\Delta\) 在 \(|\Im\theta|<\Delta/2\) 内解析。
并且由于 \(\sigma_\Delta(\theta)>0\)（\(\theta\in[-\pi,\pi]\)）且解析函数连续，任取 \(0<\varepsilon<\Delta/2\)，在闭带 \(|\Im\theta|\le \Delta/2-\varepsilon\) 上有一致下界 \(\inf|\sigma_\Delta|>0\)。
因此 \(1/\sigma_\Delta\) 在该闭带内解析且有界，其 Fourier 系数以速率 \(e^{-(\Delta/2-\varepsilon)|n|}\) 衰减。
对无限 Toeplitz 算子，逆核的 Toeplitz 系数即为 \(1/\sigma_\Delta\) 的 Fourier 系数；有限截断 \(G_\kappa(\Delta)^{-1}\) 的矩阵元衰减由同一 Wiener 代数闭包与有限截断稳定性推出。证毕。
\end{proof}

\begin{proposition}[解析带宽 \texorpdfstring{$\Delta/2$}{Delta/2} 的尖锐性与指数门槛不可超越]\label{prop:xi-hellinger-toeplitz-bandwidth-sharpness}
在命题 \ref{prop:xi-hellinger-toeplitz-symbol-poisson} 的等距设定下，定义解析带宽
\[
a_{\max}:=\sup\Bigl\{a>0:\ \sigma_\Delta\ \text{在}\ |\Im\theta|<a\ \text{内解析}\Bigr\}.
\]
则 \(a_{\max}=\Delta/2\)。
并且不存在 \(\delta>0\) 使得对所有 \(\kappa\) 与所有 \(1\le j,k\le \kappa\) 有一致界
\[
\bigl|\bigl(G_\kappa(\Delta)^{-1}\bigr)_{jk}\bigr|\ \le\ C\,e^{-(\Delta/2+\delta)|j-k|}.
\]
\end{proposition}
\begin{proof}
由 Poisson 周期和表示
\(
\sigma_\Delta(\theta)=\Delta^{-1}\sum_{k\in\ZZ}\pi^2\cosh^{-2}(\pi(\theta+2\pi k)/\Delta)
\)
与 \(\cosh(u)=0\iff u=\mathrm{i}\pi(\tfrac12+n)\)（\(n\in\ZZ\)）可知，\(\sigma_\Delta\) 在
\(
\theta=-2\pi k+\mathrm{i}\Delta(\tfrac12+n)
\)
处具有二阶极点。
因此在任意 \(a>\Delta/2\) 的条带 \(|\Im\theta|<a\) 内必包含极点，\(\sigma_\Delta\) 不可能解析；结合命题 \ref{prop:xi-hellinger-toeplitz-inverse-localization} 中 \(|\Im\theta|<\Delta/2\) 的解析性，得到 \(a_{\max}=\Delta/2\)。
\par
若存在 \(\delta>0\) 与常数 \(C\) 使所示一致界成立，则在 \(\kappa\to\infty\) 的 Toeplitz 有限截断稳定口径下，双边 Toeplitz 卷积算子 \(T_\Delta\) 的逆核系数将以速率 \(e^{-(\Delta/2+\delta)|n|}\) 绝对可和，从而 \(1/\sigma_\Delta\) 属于带权 Wiener 代数
\(
\sum_{n\in\ZZ}\abs{\widehat{(1/\sigma_\Delta)}(n)}\,e^{(\Delta/2+\delta)|n|}<\infty
\)。
由 Wiener 引理，对应的逆元 \(\sigma_\Delta\) 亦必须属于同一带权代数。
但命题 \ref{prop:xi-hellinger-toeplitz-symbol-poisson} 给出 \(\sigma_\Delta(\theta)=\sum_{n\in\ZZ}A(n\Delta)e^{\mathrm{i}n\theta}\)，而 \(A(n\Delta)\sim n\Delta\,e^{-n\Delta/2}\)（\(n\to\infty\)）使得
\(
\sum_{n\ge 1}A(n\Delta)\,e^{(\Delta/2+\delta)n}=+\infty
\)，矛盾。
证毕。
\end{proof}

\begin{theorem}[Toeplitz 热力学极限：等距体积熵密度由单一符号积分封口]\label{thm:xi-hellinger-toeplitz-szego-thermo-limit}
在命题 \ref{prop:xi-hellinger-toeplitz-symbol-poisson} 的设定下，令
\(
D_\kappa(\Delta):=\det(G_{jk})_{1\le j,k\le\kappa}
\)。
则存在极限
\[
\boxed{\;
\lim_{\kappa\to\infty}\frac{1}{\kappa}\log D_\kappa(\Delta)
=\frac{1}{2\pi}\int_{-\pi}^{\pi}\log\sigma_\Delta(\theta)\,\dd\theta\;}
\]
并因此
\[
\boxed{\;
\lim_{\kappa\to\infty}\frac{1}{\kappa}S_{1/2}(\boldsymbol{\sigma})
=-\frac{1}{2\pi}\int_{-\pi}^{\pi}\log\sigma_\Delta(\theta)\,\dd\theta\;}
\qquad(\sigma_j=\sigma_0+(j-1)\Delta).
\]
\end{theorem}
\begin{proof}
由命题 \ref{prop:xi-hellinger-toeplitz-symbol-poisson}，\(\sigma_\Delta\) 在 \([-\pi,\pi]\) 上严格正且解析。
因此 Szeg\H{o} 定理适用，给出 Toeplitz 行列式的线性极限式；第二式由 \(S_{1/2}=-\log D_\kappa\) 直接推出。证毕。
\end{proof}

\begin{proposition}[Heine--Szeg\H{o} 幺正矩阵模型：等距深度体积即 CUE 外势配分函数]\label{prop:xi-hellinger-toeplitz-heine-szego-cue}
在命题 \ref{prop:xi-hellinger-toeplitz-symbol-poisson} 的设定下，有幺正积分表示
\[
\boxed{\;
D_\kappa(\Delta)
=\frac{1}{(2\pi)^\kappa\kappa!}
\int_{[-\pi,\pi]^\kappa}
\Bigl(\prod_{r=1}^{\kappa}\sigma_\Delta(\theta_r)\Bigr)
\prod_{1\le r<\ell\le\kappa}\bigl|e^{\mathrm{i}\theta_r}-e^{\mathrm{i}\theta_\ell}\bigr|^2\,
\dd\theta_1\cdots\dd\theta_\kappa.\;}
\]
从而等距构型下的 \(S_{1/2}\) 等价于 CUE（\(\beta=2\)）模型在外势 \(-\log\sigma_\Delta\) 下的自由能增量。
\end{proposition}
\begin{proof}
这是 Toeplitz 行列式与幺正本征角积分之间的 Heine--Szeg\H{o} 恒等式。证毕。
\end{proof}

\input{para__xi-hellinger-toeplitz-tail-asymptotics}


\begin{proposition}[logistic 噪声下的多类别最优判决：最近邻规则与显式平均错误率]\label{prop:xi-logistic-multiclass-map-error}
令 \(\kappa\ge 2\)，并取有序深度对数谱
\(
\sigma_1<\sigma_2<\cdots<\sigma_\kappa
\)。
令离散输入 \(J\sim\mathrm{Unif}\{1,\dots,\kappa\}\)，观测信道为
\[
S=\sigma_J+U,
\qquad
U\sim\phi.
\]
则在等先验下最优 MAP（等价于最大似然）判决为最近邻规则
\[
\widehat J(S)\in\arg\min_{1\le j\le\kappa}|S-\sigma_j|,
\]
其决策边界为相邻中点 \((\sigma_j+\sigma_{j+1})/2\)。
并且等先验平均 Bayes 最小错误率具有严格闭式
\[
\boxed{\;
P_{\mathrm e}^{(\kappa)}
=\mathbb P(\widehat J(S)\neq J)
=\frac{2}{\kappa}\sum_{j=1}^{\kappa-1}\frac{1}{1+\exp\!\bigl((\sigma_{j+1}-\sigma_j)/2\bigr)}\; }.
\]
因此多类别可分性完全由相邻间隔 \(\Delta_j:=\sigma_{j+1}-\sigma_j\) 的 logistic 尾项决定。
\end{proposition}
\begin{proof}
因为 \(\phi\) 为偶函数且在 \([0,\infty)\) 上严格递减，故对固定 \(s\) 最大化 \(j\mapsto \phi(s-\sigma_j)\) 等价于最小化 \(j\mapsto |s-\sigma_j|\)，从而得到最近邻判决与中点边界。
\par
令 \(\Delta_j=\sigma_{j+1}-\sigma_j\)。
当 \(J=j\)（\(1<j<\kappa\)）时，误判事件等价于
\(
U\ge \Delta_j/2
\)
或
\(
U\le -\Delta_{j-1}/2
\)；
端点类 \(j=1,\kappa\) 仅有一侧尾部。
对称性给出
\(
\mathbb P(U\le -x)=\mathbb P(U\ge x)=1/(1+e^{x})
\)（\(x\ge 0\)），
对 \(j\) 求平均并将每个间隔的尾部贡献计数两次，即得所示闭式。证毕。
\end{proof}

% (User input window)

\begin{corollary}[由局部边界质量可逆重建深度间隔：错误剖面与间隔剖面一一对应]\label{cor:xi-logistic-boundary-mass-inverts-spacing}
\leanverified{paper\_xi\_logistic\_boundary\_mass\_inverts\_spacing}
在命题 \ref{prop:xi-logistic-multiclass-map-error} 的设定下，令相邻间隔 \(\Delta_j:=\sigma_{j+1}-\sigma_j>0\)。
定义局部边界质量（单侧尾概率）
\[
e_j:=\mathbb P\bigl(U\ge \Delta_j/2\bigr)=\frac{1}{1+e^{\Delta_j/2}}\in(0,1/2).
\]
则映射 \(\Delta_j\mapsto e_j\) 严格单调且显式可逆：
\[
\boxed{\;
\Delta_j=2\log\Bigl(\frac{1-e_j}{e_j}\Bigr).\;}
\]
因此一旦可观测到每个相邻切割的尾质量 \((e_j)_{j=1}^{\kappa-1}\)，深度间隔序列 \((\Delta_j)\) 可逐点无损恢复；并且
\(
P_{\mathrm e}^{(\kappa)}=\frac{2}{\kappa}\sum_{j=1}^{\kappa-1}e_j
\)
给出总错误率作为局部边界质量的线性汇总。
\end{corollary}
\begin{proof}
由 \(\mathbb P(U\ge x)=1/(1+e^{x})\)（\(x\ge 0\)）直接得到 \(e_j=(1+e^{\Delta_j/2})^{-1}\)，并由代数变形得到反演公式。
最后一式是命题 \ref{prop:xi-logistic-multiclass-map-error} 的误差闭式在记号 \(e_j\) 下的重写。证毕。
\end{proof}

\begin{corollary}[固定跨度下的最优深度布局：等距为误差极小解与必要跨度律]\label{cor:xi-logistic-optimal-spacing-span-law}
\leanverified{paper\_xi\_logistic\_optimal\_spacing\_span\_law}
在命题 \ref{prop:xi-logistic-multiclass-map-error} 的设定下，令总跨度
\[
L:=\sigma_\kappa-\sigma_1=\sum_{j=1}^{\kappa-1}\Delta_j.
\]
则有严格下界
\[
\boxed{\;
P_{\mathrm e}^{(\kappa)}
\ \ge\
\frac{2(\kappa-1)}{\kappa}\cdot\frac{1}{1+\exp\!\bigl(L/(2(\kappa-1))\bigr)}\; }.
\]
并且等号成立当且仅当
\(
\Delta_1=\cdots=\Delta_{\kappa-1}=L/(\kappa-1)
\)，即深度对数谱等距排列为唯一的误差极小解。
从而若要求 \(P_{\mathrm e}^{(\kappa)}\le \varepsilon\in(0,1)\)，则必有
\[
\boxed{\;
L\ \ge\ 2(\kappa-1)\log\!\Bigl(\frac{2(\kappa-1)}{\kappa\varepsilon}-1\Bigr)
=2(\kappa-1)\log\frac{1}{\varepsilon}+O(\kappa)\qquad(\varepsilon\downarrow0).\;}
\]
\end{corollary}
\begin{proof}
由命题 \ref{prop:xi-logistic-multiclass-map-error}，
\(
P_{\mathrm e}^{(\kappa)}=\frac{2}{\kappa}\sum_{j=1}^{\kappa-1}g(\Delta_j)
\)
其中 \(g(x)=(1+e^{x/2})^{-1}\)。
对 \(x>0\)，函数 \(g\) 可写为 \(g(x)=\bigl(1+e^{x/2}\bigr)^{-1}\) 并满足 \(g''(x)>0\)，故严格凸。
由 Jensen 不等式得到
\(
\frac{1}{\kappa-1}\sum_{j}g(\Delta_j)\ge g(L/(\kappa-1))
\)，
从而得到误差下界；严格性与等号条件给出唯一性。
跨度律由下界反解 \(L\) 得到。证毕。
\end{proof}

\begin{proposition}[互信息的熵差表示与低分辨率二次律：常数 \texorpdfstring{$1/6$}{1/6} 的刚性出现]\label{prop:xi-logistic-mutual-information-low-resolution}
在命题 \ref{prop:xi-logistic-multiclass-map-error} 的设定下，令输出混合密度
\[
p(s)=\frac{1}{\kappa}\sum_{j=1}^{\kappa}\phi(s-\sigma_j),
\]
并以自然对数定义微分熵 \(h(f):=-\int_{\RR}f\log f\)。
则互信息满足严格恒等式
\[
\boxed{\;I(J;S)=h(p)-h(\phi)=h(p)-2.\;}
\]
并且在低分辨率 regime（例如令中心化输入 \(X:=\sigma_J-\frac{1}{\kappa}\sum_{j}\sigma_j\) 且 \(\mathbb E[X^{2}]\) 足够小）存在二次近似律
\[
\boxed{\;
I(J;S)=\frac{1}{6}\,\Var(\sigma_J)+O\!\bigl(\mathbb E[|X|^{4}]\bigr).\;}
\]
其中系数 \(1/6\) 由核 \(\phi\) 的位置 Fisher 信息常数 \(\mathcal I_{\mathrm{loc}}(\phi)=1/3\) 强制确定，独立于 \(\kappa\) 与具体构型。
\end{proposition}
\begin{proof}
加性噪声信道给出 \(I(J;S)=h(S)-h(U)=h(p)-h(\phi)\)。
对 \(h(\phi)\)，作代换 \(x=e^{s}\) 得 \(\phi(s)\dd s=\dd x/(1+x)^{2}\) 且
\(
\log\phi(s)=\log x-2\log(1+x)
\)；
由 \(x\mapsto 1/x\) 的对称性可得
\(
\int_{0}^{\infty}\frac{\log x}{(1+x)^{2}}\dd x=0
\)，并且
\(
\int_{0}^{\infty}\frac{\log(1+x)}{(1+x)^{2}}\dd x=1
\)，故 \(h(\phi)=2\)。
\par
二次近似律属于加性噪声信道的小扰动展开：当输入尺度趋于零时，
\(
I(J;S)=\frac{\mathcal I_{\mathrm{loc}}(\phi)}{2}\Var(\sigma_J)+O(\mathbb E[|X|^4])
\)。
对 logistic 核，\(\phi'(u)/\phi(u)=-\tanh(u/2)\)，故
\(
\mathcal I_{\mathrm{loc}}(\phi)=\mathbb E[\tanh^{2}(U/2)]
\)。
令 \(T=\tanh(U/2)\in(-1,1)\)，则代换 \(u=2\,\operatorname{artanh}(t)\) 给出
\[
\phi(u)\,\dd u=\frac{1}{4\cosh^{2}(u/2)}\,\dd u=\frac12\,\dd t,
\]
从而 \(T\sim\mathrm{Unif}(-1,1)\) 且 \(\mathcal I_{\mathrm{loc}}(\phi)=\mathbb E[T^{2}]=1/3\)，得到系数 \(1/6\)。证毕。
\end{proof}

\leanverified{paper\_xi\_logistic\_fano\_lower\_bound\_from\_error}
\begin{corollary}[由显式平均错误率导出的 Fano 型可计算下界]\label{cor:xi-logistic-fano-lower-bound-from-error}
在命题 \ref{prop:xi-logistic-multiclass-map-error} 的设定下，有
\[
I(J;S)\ \ge\ \log\kappa\ -\ h_{\mathrm b}\!\bigl(P_{\mathrm e}^{(\kappa)}\bigr)\ -\ P_{\mathrm e}^{(\kappa)}\log(\kappa-1),
\]
其中 \(h_{\mathrm b}(u)=-u\log u-(1-u)\log(1-u)\) 为二元熵函数。
\end{corollary}
\begin{proof}
由 Fano 不等式。证毕。
\end{proof}

\begin{proposition}[位置 Fisher 信息的严格亏损：由后验方差给出的深度分散度证书]\label{prop:xi-logistic-fisher-information-deficit}
仍取
\(
p(s)=\frac{1}{\kappa}\sum_{j=1}^{\kappa}\phi(s-\sigma_j)
\)，并定义位置族 \(p_\theta(s)=p(s-\theta)\) 的 Fisher 信息
\[
\mathcal I_{\mathrm{loc}}(p):=\int_{\RR}\frac{(p'(s))^{2}}{p(s)}\,\dd s.
\]
令 \(J\sim\mathrm{Unif}\{1,\dots,\kappa\}\)、\(S=\sigma_J+U\)（\(U\sim\phi\)），并记
\[
T:=\tanh\!\Bigl(\frac{S-\sigma_J}{2}\Bigr),
\qquad
m(S):=\mathbb E[T\mid S].
\]
则存在严格恒等式
\[
\boxed{\;
\mathcal I_{\mathrm{loc}}(p)=\mathbb E[m(S)^{2}]
=\frac{1}{3}-\mathbb E\!\bigl[\Var(T\mid S)\bigr]
\ \le\ \frac{1}{3}\; }.
\]
并且等号成立当且仅当 \(\sigma_1=\cdots=\sigma_\kappa\)。
因此 \( \frac{1}{3}-\mathcal I_{\mathrm{loc}}(p)\) 是尺度曲率混合所引入的不可逆信息亏损，并以条件方差形式给出无拟合的可审计分散度证书。
\end{proposition}
\begin{proof}
由 \(\phi'(x)=-\tanh(x/2)\phi(x)\) 得
\[
p'(s)=-\frac{1}{\kappa}\sum_{j=1}^{\kappa}\tanh\!\Bigl(\frac{s-\sigma_j}{2}\Bigr)\phi(s-\sigma_j).
\]
因此
\(
-p'(s)/p(s)=\mathbb E[T\mid S=s]=m(s)
\)，从而
\(
\mathcal I_{\mathrm{loc}}(p)=\mathbb E[m(S)^{2}]
\)。
又因 \(T=\tanh(U/2)\)，同上代换 \(u=2\,\operatorname{artanh}(t)\) 得 \(T\sim\mathrm{Unif}(-1,1)\)，故 \(\mathbb E[T^{2}]=1/3\)（亦即 \(\mathcal I_{\mathrm{loc}}(\phi)=1/3\)），
由全方差分解
\(
\mathbb E[T^{2}]=\mathbb E[m(S)^{2}]+\mathbb E[\Var(T\mid S)]
\)
得到所示恒等式与上界。
若等号成立，则 \(\Var(T\mid S)=0\) 几乎处处，迫使 \(T\) 为 \(S\) 的确定函数；
在均匀先验的有限混合下，这等价于全部平移参数相同。证毕。
\end{proof}

\begin{remark}[信息几何的一维坍缩：以 \(\mathrm{TV}\) 为坐标的初等参数化]\label{rem:xi-logistic-tv-coordinate-parametrization}
令 \(v:=\|\phi_a-\phi_b\|_{\mathrm{TV}}\in[0,1)\)。
由命题 \ref{prop:xi-logistic-divergence-dictionary} 有严格反演
\[
|\Delta|=4\,\operatorname{artanh}(v)=2\log\!\Bigl(\frac{1+v}{1-v}\Bigr).
\]
并且
\[
H^{2}(\phi_a,\phi_b)
=1-\frac{(1-v^{2})\,\operatorname{artanh}(v)}{v},
\qquad
D_{\mathrm{KL}}(\phi_a\|\phi_b)
=2\frac{1+v^{2}}{v}\operatorname{artanh}(v)-2,
\]
\[
\chi^{2}(\phi_a\|\phi_b)
=\frac{16}{3}\frac{v^{2}}{(1-v^{2})^{2}},
\qquad
D_{2}(\phi_a\|\phi_b)
=\log\!\bigl(1+\chi^{2}(\phi_a\|\phi_b)\bigr).
\]
因此该平移族的若干核心散度可由单一标量坐标 \(v\) 完全参数化。
\end{remark}

\begin{remark}[Fisher--Rao 度量的显式化与二次近似偏离函数]\label{rem:xi-logistic-fisher-rao-and-curvature}
尺度核平移族 \(\{\phi_a\}\) 的位置 Fisher 信息为常数 \(\mathcal I_{\mathrm{loc}}(\phi)=1/3\)，故 Fisher--Rao 度量在参数 \(a\) 上为常系数，
任意两点的 Fisher--Rao 距离满足
\[
d_{\mathrm{FR}}(\phi_a,\phi_b)=\sqrt{\mathcal I_{\mathrm{loc}}(\phi)}\,|a-b|=\frac{|\Delta|}{\sqrt3}.
\]
并且由命题 \ref{prop:xi-logistic-divergence-dictionary} 的闭式可定义偏离函数
\[
\Psi(\Delta):=D_{\mathrm{KL}}(\phi_a\|\phi_b)-\frac{1}{2}d_{\mathrm{FR}}(\phi_a,\phi_b)^{2}
=|\Delta|\coth(|\Delta|/2)-2-\frac{\Delta^{2}}{6},
\]
其在 \(\Delta\to 0\) 时满足 \(\Psi(\Delta)=O(\Delta^{4})\)，并刻画 KL 与 Fisher--Rao 二次型在非线性可见尺度上的偏离。
\end{remark}

\begin{corollary}[多深度判别的边界质量表示：平均错误率的相邻切割分解]\label{cor:xi-logistic-boundary-mass-error}
在命题 \ref{prop:xi-logistic-multiclass-map-error} 的设定下，令 \(\Delta_j=\sigma_{j+1}-\sigma_j\)。
则平均错误率具有严格边界质量形式
\[
P_{\mathrm e}^{(\kappa)}
=\frac{2}{\kappa}\sum_{j=1}^{\kappa-1}\mathbb P\!\Bigl(U\ge \frac{\Delta_j}{2}\Bigr),
\qquad
\mathbb P(U\ge x)=\frac{1}{1+e^{x}}\quad(x\ge 0).
\]
因此 logistic 尺度核下的多类别误差不存在高维干涉项，仅由相邻间隔诱导的边界尾质量线性叠加给出。
\end{corollary}
\begin{proof}
由命题 \ref{prop:xi-logistic-multiclass-map-error} 的推导，误判等价于噪声落入相邻中点切割的两侧外尾，并由对称性化为单侧尾概率之和。证毕。
\end{proof}

\begin{remark}[\texorpdfstring{$4$}{4} 阶闭合阈值的散度表征]\label{rem:xi-logistic-fourth-order-closure}
命题 \ref{prop:xi-logistic-divergence-dictionary} 的小间隔展开显示：\(\mathrm{TV}\) 首项为线性（系数 \(1/4\)），而 \(H^{2}\) 与 \(D_{\mathrm{KL}}\) 的首个非平凡项为二次；
并且 \(D_{\mathrm{KL}}=4H^{2}+O(\Delta^{4})\) 把二者在四阶精度内锁定为同一局部几何。
因此当同时关注边界切割（由 \(\mathrm{TV}\) 的线性尾控制）与熵势涨落（由 \(H^{2}\)、\(D_{\mathrm{KL}}\) 的二次项控制）时，四阶项成为最小的闭合阶。
\end{remark}

\begin{corollary}[以信息可见性为目标的深度谱设计律：最优构型收缩到等距族]\label{cor:xi-logistic-depth-design-law}
在固定 \(\kappa\) 与总跨度 \(L\) 的约束下，使命题 \ref{prop:xi-logistic-multiclass-map-error} 的平均错误率 \(P_{\mathrm e}^{(\kappa)}\) 最小的构型唯一地由等距谱实现：
\[
\sigma_j=\sigma_1+\frac{j-1}{\kappa-1}L\qquad(1\le j\le\kappa).
\]
并且该构型同时使推论 \ref{cor:xi-logistic-fano-lower-bound-from-error} 给出的互信息下界达到最大值。
\end{corollary}
\begin{proof}
由推论 \ref{cor:xi-logistic-optimal-spacing-span-law}，在固定跨度下 \(P_{\mathrm e}^{(\kappa)}\) 的下界由 Jensen 取等条件唯一实现；
而 Fano 下界对 \(P_{\mathrm e}^{(\kappa)}\) 单调递减，故同一构型也最大化该下界。证毕。
\end{proof}



\begin{proposition}[截断 Pad\'e--Stieltjes 逼近的保序性：极点实负、留数正、并给出全尺度单调收敛]\label{prop:xi-pade-stieltjes-order-preserving}
\leanverified{paper\_xi\_pade\_stieltjes\_order\_preserving}
若 \(S\) 为上述 Stieltjes 函数，则其在 \(t=0\) 处的对角 Pad\'e 逼近 \([N/N]\)（匹配前 \(2N\) 个 Taylor 系数）仍为 Stieltjes 型有理函数：其全部极点位于 \((-\infty,-1)\) 且为单根，并且对应部分分式展开的留数严格为正。
此外，对任意 \(t>0\)，序列 \([N/N](t)\) 在 \(N\) 上按奇偶分裂为两条单调链，分别从上、下夹逼真值 \(S(t)\)；并且当 \(N\to\infty\) 时其夹逼区间收敛到单点 \(\{S(t)\}\)。
该结论把“有限寄存器证书”的可审计性强化为保序性：逼近过程中既不产生伪极点，也不产生符号翻转的伪权重，从而在制度层面排除了由近似引入的非物理缺陷指数。
\end{proposition}
\begin{proof}[证明要点]
Stieltjes 类的 Pad\'e 逼近保留正权重与实负极点结构，是连分式/正交多项式理论的直接推论；其奇偶单调性对应 Stieltjes 连分式收敛子的交错夹逼性质，而收敛性来自 Stieltjes 变换在 \((0,\infty)\) 上的完全单调与紧性控制。证毕。
\end{proof}

\begin{proposition}[深度最浅原子的纯尾部可辨识：矩比极限决定 \texorpdfstring{$a_{\min}$}{a_min} 并给出指数收敛率]\label{prop:xi-moment-ratio-identifies-amin}
\leanverified{paper\_xi\_moment\_ratio\_identifies\_amin}
在有限原子口径下，设
\(
\mu=\sum_{j=1}^{\kappa} w_j\,\delta_{a_j}
\)
且 \(1<a_1:=a_{\min}<a_2\le\cdots\le a_{\kappa}\)、\(w_j>0\)。
令矩链
\[
\mu_n:=\int a^{-(n+1)}\,\dd\mu(a)=\sum_{j=1}^{\kappa} w_j\,a_j^{-(n+1)}\qquad(n\ge 0).
\]
则存在极限
\[
\boxed{\;
\lim_{n\to\infty}\frac{\mu_n}{\mu_{n+1}}=a_1,
\qquad
\lim_{n\to\infty}\frac{\mu_{n+1}}{\mu_n}=\frac{1}{a_1}\; }.
\]
并且收敛为指数级：存在常数 \(C>0\) 使
\[
\boxed{\;
\left|\frac{\mu_n}{\mu_{n+1}}-a_1\right|
\ \le\ C\left(\frac{a_1}{a_2}\right)^{n}\qquad(n\to\infty)\; }.
\]
因此最浅深度 \(\delta_{\min}=a_1-1\) 在不引入任何相位载体且不做全谱反演的条件下，可由单变量矩链尾部的比值极限唯一恢复；谱隙比 \(a_1/a_2\) 精确控制其制度内可辨识速度。
\end{proposition}
\begin{proof}
把矩链写作 \(\mu_n=w_1a_1^{-(n+1)}(1+R_n)\)，其中
\(
R_n:=\sum_{j=2}^{\kappa}\frac{w_j}{w_1}\bigl(\frac{a_1}{a_j}\bigr)^{n+1}
\)。
则 \(|R_n|\le C_0(\frac{a_1}{a_2})^{n}\)（某常数 \(C_0>0\)），并且
\[
\frac{\mu_n}{\mu_{n+1}}
=a_1\cdot\frac{1+R_n}{1+R_{n+1}}
=a_1\left(1+\frac{R_n-R_{n+1}}{1+R_{n+1}}\right).
\]
由 \(R_n-R_{n+1}=O((a_1/a_2)^n)\) 得所示指数误差界。证毕。
\end{proof}

\begin{proposition}[Laplace 配分尾率的极限律：\texorpdfstring{$\delta_{\min}$}{delta_min} 为配分函数的速率]\label{prop:xi-partition-tail-rate-deltamin}
\leanverified{paper\_xi\_partition\_tail\_rate\_deltamin}
在上一命题的有限原子设定下，令 \(\delta_j:=a_j-1\)，并定义深度配分函数
\[
Z_\delta(u):=\sum_{j=1}^{\kappa} w_j\,e^{-u\delta_j}\qquad(u\ge 0).
\]
则存在极限
\[
\boxed{\;
\lim_{u\to\infty}-\frac{1}{u}\log Z_\delta(u)=\delta_{\min}:=\min_j \delta_j\; }.
\]
并且若 \(\delta_2\) 为第二小深度，则有指数误差界：存在常数 \(C>0\) 使
\[
\left|-\frac{1}{u}\log Z_\delta(u)-\delta_{\min}\right|
\ \le\ \frac{1}{u}\log\!\left(1+C\,e^{-u(\delta_2-\delta_{\min})}\right).
\]
该结论把“最浅深度”的可辨识性提升为热力学极限律：\(\delta_{\min}\) 为配分函数的速率函数，而深度谱隙 \(\delta_2-\delta_{\min}\) 决定其收敛的指数尺度。
\end{proposition}
\begin{proof}
写作
\(
Z_\delta(u)=e^{-u\delta_{\min}}\bigl(w_1+\sum_{j\ge 2} w_j e^{-u(\delta_j-\delta_{\min})}\bigr)
\)，
取对数并除以 \(u\) 即得极限。
误差界由括号内第二项的指数衰减给出。证毕。
\end{proof}

\begin{corollary}[尾率证书的最小观测尺度：\texorpdfstring{$\delta_{\min}$}{delta_min} 驱动的不可区分下界]\label{cor:xi-partition-tail-min-scale-lower-bound}
\leanverified{paper\_xi\_partition\_tail\_min\_scale\_lower\_bound}
仍在命题 \ref{prop:xi-partition-tail-rate-deltamin} 的有限缺陷设定下，设制度仅允许在离散网格
\(
u\in\{0,\Delta u,2\Delta u,\dots,M\Delta u\}
\)
上观测配分尾率 \(Z_\delta(u)\)（或等价观测 \(-\frac{1}{u}\log Z_\delta(u)\)），且观测精度要求达到相对误差 \(\varepsilon\in(0,1)\) 的量级。
则要在最坏情形上把“临界无缺陷（\(\delta_{\min}=0\)，在闭包意义下）”与“存在最浅深度 \(\delta_{\min}>0\)”区分开来，所需最大观测尺度必满足
\[
\boxed{\;
M\Delta u\ \gtrsim\ \frac{1}{\delta_{\min}}\log\!\frac{1}{\varepsilon}\;}
\]
其中隐含常数仅依赖于权重比的先验界（例如 \(w_2/w_1\) 的上界）。
该下界刻画了制度内的不可见边界层：\(\delta_{\min}\) 越小，必须把观测推进到越大的 \(u\)（更长的“制度时间”）才能产生可审计分离。
\end{corollary}
\begin{proof}[证明要点]
由 \(Z_\delta(u)\ge w_1e^{-u\delta_{\min}}\) 与 \(Z_\delta(u)\le (w_1+\cdots+w_\kappa)e^{-u\delta_{\min}}\) 得到：若 \(\delta_{\min}=0\)（闭包口径）则 \(Z_\delta(u)\ge w_1\) 不随 \(u\) 衰减；而若 \(\delta_{\min}>0\) 则 \(Z_\delta(u)\le \Phi\,e^{-u\delta_{\min}}\) 指数衰减（\(\Phi:=\sum_j w_j\)）。
因此要把二者在相对误差 \(\varepsilon\) 下区分，必须进入指数尾部区间 \(e^{-u\delta_{\min}}\lesssim\varepsilon\)，即 \(u\gtrsim \delta_{\min}^{-1}\log(1/\varepsilon)\)。
若 \(u\) 仅能取到 \(M\Delta u\)，即得所示下界。证毕。
\end{proof}

\endinput
